%%%%%%%%%%%%%%%%%%%%%%%%%%%%%%%%%%%%%%%%%%%%%%%%%%%%%%%%%%%%%%%%%%%%%%
% Lớp 12 | Tọa độ của véc-tơ trong không gian | Tọa độ của véc-tơ và các công thức
% Dạng 2H2?2-2: Tìm tọa độ điểm
% Số câu: 79  (sinh tự động bởi tach_cau.py; từ nay thầy sửa trực tiếp file này)
%%%%%%%%%%%%%%%%%%%%%%%%%%%%%%%%%%%%%%%%%%%%%%%%%%%%%%%%%%%%%%%%%%%%%%

%%%%%%%%%%%%%%%% TRẮC NGHIỆM 4 PHƯƠNG ÁN %%%%%%%%%%%%%%%%

% nguồn: Lop_12_Hinh_hoc_Chuong_2
\begin{ex}%[2H2N2-2]%[ma:1c1f3067]
	Trong không gian $Oxyz$, hình chiếu vuông  góc của điểm $A(1;2;5)$ trên mặt $Oxz$ có tọa độ là
	\choice 
		 {$(0;2;5)$}
		 {$(0;2;0)$}
		{\True $(1;0;5)$} 
		 {$(0;0;5)$}
	\loigiai{ 
		Hình chiếu vuông góc của điểm  $A(1;2;5)$ trên mặt $Oxz$ có tọa độ là $(1;0;5)$.
	 }
\end{ex}

% nguồn: Lop_12_Hinh_hoc_Chuong_2
\begin{ex}%[2H2N2-2]%[ma:d87b8861]
	Trong không gian $Oxyz$, hình chiếu vuông góc điểm $A(2;1;3)$ lên mặt phẳng $(Oxy)$ là 
	 \choice
		 {$(-2;-1;3)$}
		 {\True $(2;1;0)$}
		{$(-2;1;3)$} 
		{$(-2;-1;-3)$} 
	 \loigiai{
		Hình chiếu vuông  góc của điểm $A(2;1;3)$ lên mặt phẳng $(Oxy)$ là điểm $(2;1;0)$.
	 }
\end{ex}

% nguồn: Lop_12_Hinh_hoc_Chuong_2
\begin{ex}%[2H2N2-2]%[ma:72d6496d]
	Trong không gian $Oxyz$, cho điểm $A(1;-2;3)$. Hình chiếu vuông góc của điểm $A$ lên mặt phẳng $(Oyz)$  có tọa độ là
	\choice 
		{$(1;0;3)$} 
		 {$(1;0;0)$}
		 {$(1;-2;0)$}
		{\True  $(0;-2;3)$}
	 \loigiai{
		Để tìm tọa  độ hình chiếu của điểm $A(1;-2;3)$ lên mặt phẳng $(Oyz)$ ta chỉ cần giữ nguyên tung độ và cao độ, cho hoành độ bằng $0$.
	 }
\end{ex}

% nguồn: Lop_12_Hinh_hoc_Chuong_2
\begin{ex}%[2H2N2-2]%[ma:a956be2b]
	Trong không gian với hệ tọa độ $Oxyz$, tọa độ  hình chiếu của điểm $A(-2;-1;3)$ trên mặt phẳng $Oyz$ là
	 \choice
		{$(0;-1;0)$} 
		{$(-2;0;0)$} 
		{\True $(0;-1;3)$} 
		 {$(-2;-1;0)$}
	 \loigiai{
		Tọa độ hình chiếu  của điểm $A(-2;-1;3)$ trên mặt phẳng $Oyz$ là $(0;-1;3)$.
	} 
\end{ex}

% nguồn: Lop_12_Hinh_hoc_Chuong_2
\begin{ex}%[2H2N2-2]%[ma:bab25204]
	Trong không gian với hệ tọa độ $Oxyz$, hình chiếu vuông  góc điểm $A(-1;2;3)$ lên mặt phẳng $(Oxy)$ là
	 \choice
		{$(0;0;3)$} 
		{\True$(-1;2;0)$} 
		 {$(-1;2;3)$}
		 {$(1;-2;-3)$}
	 \loigiai{
		 Hình chiếu vuông góc điểm $A(-1;2;3)$ lên mặt phẳng $(Oxy)$ là điểm $A'(-1;2;0)$.
	} 
\end{ex}

% nguồn: Lop_12_Hinh_hoc_Chuong_2
\begin{ex}%[2H2N2-2]%[ma:db567c54]
	Trong không gian $Oxyz$, cho điểm $M$ thỏa mãn hệ thức $\overrightarrow{OM}=2\overrightarrow{i}+\overrightarrow{j}$. Tọa  độ điểm $M$ là
	 \choice
		 {$M=(0;2;1)$}
		 {$M=(1;2;0)$}
		{\True  $M=(2;1;0)$}
		 {$M=(2;0;1)$}
	\loigiai{ 
		Vì điểm $M$ thỏa mãn hệ  thức $\overrightarrow{OM}=2\overrightarrow{i}+\overrightarrow{j}$ nên tọa độ điểm $M=(2;1;0)$.
	 }
\end{ex}

% nguồn: Lop_12_Hinh_hoc_Chuong_2
\begin{ex}%[2H2N2-2]%[ma:9b3c6b32]
	Trong không gian $Oxyz$, cho tam giác $ABC$  với $A(1;2;-1)$, $B(2;-1;3$), $C(-3;5;1)$. Tọa độ trọng tâm của tam giác $ABC$ là
	\choice 
		{$(0;-2;1)$} 
		 {$(0;-2;-1)$}
		{$(0;2;-1)$} 
		{\True  $(0;2;1)$}
	\loigiai{ 
		Trọng tâm $G$ của tam giác $ABC$ có tọa độ  $\left(\dfrac{1+2+(-3)}{3};\dfrac{2+(-1)+5}{3};\dfrac{-1+3+1}{3}\right)$ hay G$\left(0;2;1\right)$.
	} 
\end{ex}

% nguồn: Lop_12_Hinh_hoc_Chuong_2
\begin{ex}%[2H2N2-2]%[ma:518e7c7f]
	Cho điểm  $A(3;-1;1)$. Hình chiếu vuông góc của điểm $A$ trên mặt phẳng $(Oyz)$ là điểm
	\choice 
		{$M(3;0;0)$} 
		{\True $N(0;-1;1)$} 
		 {$P(0;-1;0)$}
		{$Q(0;0;1)$} 
	 \loigiai{
		Ta có hình chiếu vuông góc của điểm $A$ trên mặt  phẳng ($Oyz$) là điểm $N(0;-1;1)$.
	} 
\end{ex}

% nguồn: Lop_12_Hinh_hoc_Chuong_2
\begin{ex}%[2H2N2-2]%[ma:e3429ed2]
	Trong không gian $Oxyz$, hình chiếu vuông góc của điểm  $A(1;2;5)$ lên trục $Ox$ có tọa độ là
	 \choice
		 {$(0;2;0)$}
		 {$(0;0;5)$}
		 {\True $(1;0;0)$}
		{$(0;2;5)$} 
	 \loigiai{
		Hình chiếu vuông góc của điểm $A(1;2;5)$  lên trục $Ox$ có tọa độ là $(1;0;0)$.
	 }
\end{ex}

% nguồn: Lop_12_Hinh_hoc_Chuong_2
\begin{ex}%[2H2N2-2]%[ma:41996883]
	Trong không  gian $Oxyz$, cho điểm $A(2;1;3)$ và điểm $B(4;-3;1)$. Tọa độ trung điểm $I$ của đoạn thẳng $AB$ là
	 \choice
		 {$(6;-2;4)$}
		 {\True $(3;-1;2)$}
		{$(1;-2;-1)$} 
		{$(2;-4;-2)$} 
	\loigiai{ 
		Tọa độ  trung điểm $I$ của đoạn thẳng $AB$ là
		$I\left(\dfrac{2+4}{2}; \dfrac{1+(-3)}{2}; \dfrac{3+1}{2}\right) = I(3;-1;2)$. 
	} 
\end{ex}

% nguồn: Lop_12_Hinh_hoc_Chuong_2
\begin{ex}%[2H2N2-2]%[ma:d23307f0]
	Trong không gian $Oxyz$, cho điểm $M$ thoả mãn $\overrightarrow{OM}=2 \vec{i}+3\vec{j}-4\vec{k}$. Toạ độ của  điểm $M$~là
	 \choice
		{$(-4;3;2)$} 
		 {\True $(2;3;-4)$}
		 {$(3;-4;2)$}
		{$(-2;-3;4)$} 
	 \loigiai{Ta có $\overrightarrow{OM}=2 \vec{i}+3\vec{j}-4\vec{k}$.\\
		Vậy toạ độ điểm $M$~là  $M(2;3;-4)$.
	} 
\end{ex}

% nguồn: Lop_12_Hinh_hoc_Chuong_2
\begin{ex}%[2H2N2-2]%[ma:aa0f9058]
	Trong không gian  $Oxyz$, cho điểm $A(1;1;1)$. Tìm tọa độ hình chiếu vuông góc của điểm $A$ trên mặt phẳng ($Oxz$).
	\choice 
		{\True  $(1;0;1)$}
		 {$(0;1;0)$}
		 {$(1;1;0)$}
		{$(0;1;1)$} 
	 \loigiai{
		Tọa độ hình chiếu vuông góc của  điểm $A$ trên mặt phẳng $(Oxz)$ là $H(1;0;1)$.
	} 
\end{ex}

% nguồn: Lop_12_Hinh_hoc_Chuong_2
\begin{ex}%[2H2N2-2]%[ma:2eb7c922]
	Trong  không gian $Oxyz$, cho điểm $G(1;-2;3)$ và ba điểm $A(a;0;0)$, $B(0;b;0)$, $C(0;0;c)$. Biết $G$ là trọng tâm của tam giác $ABC$. Khi đó, $a+b+c$ bằng
	\choice 
		 {$3$}
		{$9$} 
		{\True  $6$}
		 {$0$}
	 \loigiai{
		Vì $G$ là  trọng tâm của tam giác $ABC$ nên tọa độ của $G$ là
		$$\heva{&x_G=\dfrac{x_A+x_B+x_C}{3} \\&y_G=\dfrac{y_A+y_B+y_C}{3}\\&z_G=\dfrac{z_A+z_B+z_C}{3}}\Leftrightarrow \heva{&1=\dfrac{a+0+0}{3}\\&-2=\dfrac{0+b+0}{3}\\&3=\dfrac{0+0+c}{3} }\Leftrightarrow  \heva{&a=3\\&b=-6\\&c=9.}$$
		Khi  đó, $a+b+c=3+(-6)+9=6$.
	} 
\end{ex}

% nguồn: Lop_12_Hinh_hoc_Chuong_2
\begin{ex}%[2H2N2-2]%[ma:1e8e8c56]
	Trong không gian $Oxyz$,  cho hai điểm $A(2;1;1)$, $B(-1;2;1)$. Tọa độ trung điểm của đoạn thẳng $AB$ là
	 \choice
		{$I\left(\dfrac{1}{3};\dfrac{2}{3};1\right)$} 
		{$I(-3;1;0)$} 
		{$I\left(\dfrac{3}{2};-\dfrac{1}{2};0\right)$} 
		{\True $I\left(\dfrac{1}{2};\dfrac{3}{2};1\right)$} 
	\loigiai{ 
		Tọa độ trung điểm $I$  của đoạn thẳng $AB$ là
		$I\left(\dfrac{2+(-1)}{2};\dfrac{1+2}{2};\dfrac{1+1}{2}\right) = I\left(\dfrac{1}{2};\dfrac{3}{2};1\right)$. 
	 }
\end{ex}

% nguồn: Lop_12_Hinh_hoc_Chuong_2
\begin{ex}%[2H2N2-2]%[ma:dda39787]
	Trong  không gian $Oxyz$, cho hai điểm $A(1;-1; 2)$ và $B(-1; 5;-4)$. Trung điểm $M$ của đoạn thẳng $AB$ có tọa độ là
	\choice 
		{\True $(0; 2;-1)$} 
		 {$(-1; 3;-3)$}
		 {$(-2; 6;-6)$}
		 {$(0; 4;-2)$}
	 \loigiai{
		Trung điểm $M$ của  đoạn thẳng $AB$ có tọa độ là $\left(\dfrac{1-1}{2}; \dfrac{-1+5}{2}; \dfrac{2-4}{2}\right)=(0; 2;-1)$.
	 }
\end{ex}

% nguồn: Lop_12_Hinh_hoc_Chuong_2
\begin{ex}%[2H2H2-2]%[ma:09c872fc]
	 Trong không gian với hệ tọa độ $Oxyz$ cho hình bình hành $ABCD$ và các đỉnh có tọa độ lần lượt là $A(3;1;2)$, $B(1;0;1)$, $C(2;3;0)$. Tọa độ đỉnh $D$ là
	\choice 
		 {$D(1;1;0)$}
		{$D(0;2;-1)$} 
		{\True $D(4;4;1)$} 
		 {$D(1;3;-1)$}
	 \loigiai{
		Vì  $ABCD$ là hình bình hành suy ra $\overrightarrow{AD}=\overrightarrow{BC}\Leftrightarrow\heva{&x_D-3=1\\&y_D-1=3\\&z_D-2=-1}\Leftrightarrow\heva{&x_D=4\\&y_D=4\\&z_D=1}\Rightarrow D(4;4;1)$.
	 }
\end{ex}

% nguồn: Lop_12_Hinh_hoc_Chuong_2
\begin{ex}%[2H2H2-2]%[ma:7e0fc5dd]
	Trong không gian với hệ tọa độ $Oxyz$, cho  hai điểm $A(2;1;1)$, $B(-1;2;1)$. Tìm tọa độ của điểm $A'$ đối xứng với điểm $A$ qua điểm $B$?
	 \choice
		 {$A'(3;4;-3)$}
		{\True  $A'(-4;3;1)$}
		{$A'(1;3;2)$} 
		{$A'(5;0;1)$} 
	\loigiai{ 
		Điểm $A'$  đối xứng với điểm $A$ qua điểm $B$ nên $B$ là trung điểm của đoạn $AA'$.\\
		 Do đó $\heva{&x_{A'}=2x_B-x_A=-4\\&y_{A'}=2y_B-y_A=3\\&z_{A'}=2z_B-z_A=1}\Rightarrow A'(-4;3;1)$.
	} 
\end{ex}

% nguồn: Lop_12_Hinh_hoc_Chuong_2
\begin{ex}%[2H2H2-2]%[ma:847ba650]
	Trong không gian $Oxyz$ cho ba điểm $M(2;0;0)$, $N(0;-3;0)$, $P(0;0;4)$.  Nếu $MNPQ$ là hình bình hành thì tọa độ điểm $Q$ là
	\choice 
		{$(-2;-3;4)$} 
		{$(-2;-3;-4)$} 
		{\True $(2;3;4)$} 
		{$(3;4;2)$} 
	 \loigiai{
		Gọi  $Q(x;y;z)$ khi đó $\overrightarrow{MN}=(-2;-3;0)$, $\overrightarrow{QP}=(-x;-y;4-z)$.\\
		Tứ giác $MNPQ$ là  hình bình hành $\Leftrightarrow \overrightarrow{MN}=\overrightarrow{QP}\Leftrightarrow\heva{&-x=-2\\&-y=-3\\&4-z=0}\Leftrightarrow\heva{&x=2\\&y=3\\&z=4}$.\\
		 Vậy $Q(2;3;4)$.
	} 
\end{ex}

% nguồn: Lop_12_Hinh_hoc_Chuong_2
\begin{ex}%[2H2H2-2]%[ma:ee161dfd]
	Trong không gian $Oxyz$, cho ba điểm $A(1;2;-1)$,  $B(2;-1;3)$, $C(-3;5;1)$. Tìm tọa độ điểm $D$ sao cho tứ giác $ABCD$ là hình bình hành.
	\choice 
		{$D(-4;8;-5)$} 
		{$D(-2;2; 5)$} 
		{\True $D(-4;8;-3)$} 
		{$D(-2;8;-3)$} 
	 \loigiai{
		Đặt  $D(x;y;z)$.\\
		Tứ giác $ABCD$ là hình bình hành suy ra $\overrightarrow{AB}=\overrightarrow{DC}$, khi đó ta có 
		$$\heva{&1=-3-x\\&-3=5-y\\&4=1-z}\Rightarrow  \heva{&x=-4\\&y=8\\&z=-3}\Rightarrow D(-4;8;-3).$$
	} 
\end{ex}

% nguồn: Lop_12_Hinh_hoc_Chuong_2
\begin{ex}%[2H2H2-2]%[ma:eea4b72b]
	Trong không gian $Oxyz$, cho véctơ $\vec{u}=(4;-2;-3)$ và điểm $A(1;2;3)$. Tọa độ điểm $M$  thỏa mãn $\overrightarrow{AM}=\vec{u}$ là
	\choice 
		 {$M(-5;0;0)$}
		 {\True $M(5;0;0)$}
		{$M(5;-2;-3)$} 
		{$B(4;-2;-3)$} 
	\loigiai{ 
		Gọi $M(x;y;z)\Rightarrow \overrightarrow{AM}=(x-1;y-2;z-3)$.\\ 
		Ta có 
		 $\overrightarrow{AM}=\overrightarrow{u}\Leftrightarrow \heva{
		& x-1=4  \\
		&  y-2=-2 \\
		& z-3=-3  \\
		}\Leftrightarrow \heva{ 
		 & x=5 \\
		&  y=0 \\
		& z=0.}$\\ 
		Vậy $M(5;0;0)$. 
	 }
\end{ex}

% nguồn: Lop_12_Hinh_hoc_Chuong_2
\begin{ex}%[2H2H2-2]%[ma:3d95428b]
	Trong không gian $Oxyz$, cho tam giác $ABC$ có  các điểm $A(1;0;3)$, $B(2;3;-4)$, $C(-3;1;2)$. Tìm toạ độ điểm $D$ sao cho tứ giác $ABCD$ là hình bình hành.
	 \choice
		 {$D(4;2;9)$}
		{$D(-2;4;-5)$} 
		{\True $D(-4;-2;9)$} 
		 {$D(6;2;-3)$}
	 \loigiai{
		Gọi  $D(x;y;z) \Rightarrow\overrightarrow{CD}=(x+3;y-1;z-2)$ và $\overrightarrow{BA}=(-1;-3;7)$.\\
		Tứ giác $ABCD$ là hình bình hành khi và chỉ khi \begin{center}
		$\overrightarrow{BA}=\overrightarrow{CD}\Leftrightarrow\heva{ 
		& x+3=-1\\ 
		&  y-1=-3\\
		& z-2=7} 
		 \Leftrightarrow\heva{
		&  x=-4\\
		 & y=-2\\
		&  z=9}\Rightarrow C(-4;-2;9)$.
	\end{center}
} 
\end{ex}

% nguồn: Lop_12_Hinh_hoc_Chuong_2
\begin{ex}%[2H2H2-2]%[ma:54f135d4]
	 Trong không gian với hệ trục $Oxyz$ cho ba điểm $A(-1;2;-3)$, $B(1;0;2)$, $C(x;y;-2)$ thẳng hàng. Khi đó $x+y$ bằng
	 \choice
		 {\True $x+y=1$}
		 {$x+y=17$}
		{$x+y=-\dfrac{11}{5}$} 
		 {$x+y=\dfrac{11}{5}$}
	\loigiai{ 
		Ta có $\overrightarrow{AB}=(2;-2;5)$, $\overrightarrow{AC}=(x+1;y-2;1)$.\\ 
		Ba điểm $A$, $B$, $C$ thẳng hàng khi và chỉ khi $\overrightarrow{AB}$,  $\overrightarrow{AC}$ cùng phương.\\
		Suy  ra
		 $ \dfrac{x+1}{2}=\dfrac{y-2}{-2}=\dfrac{1}{5}\Leftrightarrow\heva{&x=-\dfrac{3}{5}\\&y=\dfrac{8}{5}}$.\\
		Vậy  $x+y=1$.
	} 
\end{ex}

% nguồn: Lop_12_Hinh_hoc_Chuong_2
\begin{ex}%[2H2H2-2]%[ma:128dae2d]
	Cho hình lập phương $ABCD.A'B'C'D'$ có cạnh bằng $2$. Với hệ toạ độ $Oxyz$ được thiết lập như  hình bên (gốc tọa độ $O$ trùng với điểm $A$), tọa độ điểm $C'$ là
	\begin{center}
		\begin{tikzpicture}[scale=0.8,font=\footnotesize,line join=round,line cap=round,>=stealth]
				\coordinate (A)  at (0,0);
				\coordinate (B) at  (-2,-2);
				\coordinate (D) at  (6,0);
				\coordinate (C) at  ($(B)+(D)-(A)$);
				\foreach  \i in {A,B,C,D}{\coordinate (\i') at ($(\i)+(0,4)$);}
				\draw (A')--(B')--(C')--(D')--cycle; 
				\draw (B)--(B') 
					(C)--(C') 
					(D)--(D') 
					 (B)--(C)--(D);
				\draw[dashed,thin](B)--(A)--(A') 
					(A)--(D); 
				 \foreach \i/\g in {A'/110,B'/90,C'/90,D'/-30,A/-90,B/-90,C/-90,D/-90}{\draw[fill=blue](\i) circle (1.5pt)
					($(\i)+(\g:5mm)$) node[scale=1]{$\i$};} 
				\coordinate  (Z) at ($(A)+(0,5)$);
				\coordinate (Y)  at ($(D)+(1,0)$);
				\coordinate (X) at  ($(B)-(1,1)$);
				 \draw[->] (A')--(Z);
				\draw[->] (B)--(X); 
				\draw[->]  (D)--(Y);
				\draw (X) node[scale=1,left]{$x$};  x
				\draw (Y)  node[scale=1,above]{$y$};
				\draw (Z)  node[scale=1,left]{$z$};
			\end{tikzpicture}
	\end{center}
	\choice 
		{$C'\left(2; 2; 0\right)$} 
		{\True $C'\left(2;  2; 2\right)$}
		 {$C'\left(2;0;0\right)$}
		 {$C'\left(2; 0; 2\right)$}
	\loigiai{ 
		Với $A$ là  gốc toạ độ thì $C'\left(2; 2; 2\right)$.
	} 
\end{ex}

% nguồn: Lop_12_Hinh_hoc_Chuong_2
\begin{ex}%[2H2H2-2]%[ma:88a8ca97]
	 Trong không gian với hệ tọa độ $Oxyz$, cho điểm $A(2;-1;0)$ và điểm $B(3;1;1)$. Tọa độ điểm đối xứng với $A$ qua $B$ là
	\choice 
		 {$(1;-2;-4)$}
		 {$(0;3;-1)$}
		 {\True $(4;3;2)$}
		 {$(0;-1;3)$}
	 \loigiai{
		Gọi $A'(x;y;z)$ là điểm đối xứng với $A$ qua  $B$, suy ra $B$ là trung điểm đoạn thẳng $AA'$.\\
		 Ta có $\heva{&3=\dfrac{2+x}{2}\\&1=\dfrac{-1+y}{2}\\&1=\dfrac{0+z}{2}}\Leftrightarrow\heva{&x=4\\&y=3\\&z=2.}$\\
		Suy  ra $A'(4;3;2)$.
	} 
\end{ex}

% nguồn: Lop_12_Hinh_hoc_Chuong_2
\begin{ex}%[2H2H2-2]%[ma:89058eeb]
	Trong không gian với hệ tọa độ $Oxyz$, cho điểm $M(2;-1;-3)$. Tìm tọa độ của điểm $M'$ đối xứng với điểm $M$  qua trục $Oy$.
	\choice 
		{$M'(2;1;-3)$} 
		{\True $M'(-2;-1;3)$} 
		{$M'(2;-1;-3)$} 
		 {$M'(-2;-1;-3)$}
	 \loigiai{
		Điểm $M\left(x_0;y_0;z_0\right)$  có điểm đối xứng qua trục $Oy$ là điểm $M'\left(-x_0;y_0;-z_0\right)$.\\
		Vậy  điểm $M(2;-1;-3)$ có điểm đối xứng qua trục $Oy$ là điểm $M'(-2;-1;3)$.
	 }
\end{ex}

% nguồn: Lop_12_Hinh_hoc_Chuong_2
\begin{ex}%[2H2H2-2]%[ma:c1295684]
	 Trong không gian với hệ toạ độ $Oxyz$, cho ba điểm $A(1;1;1)$, $B(5;-1;2)$, $C(3;2;-4)$. Tìm toạ độ điểm $M$ thoả mãn $\overrightarrow{MA}+2\overrightarrow{MB}-\overrightarrow{MC}=\overrightarrow{0}$
	 \choice
		 {$M\left(4;-\dfrac{3}{2};-\dfrac{9}{2}\right)$}
		 {$M\left(4;\dfrac{3}{2};\dfrac{9}{2}\right)$}
		 {\True $M\left(4;-\dfrac{3}{2};\dfrac{9}{2}\right)$}
		 {$M\left(-4;-\dfrac{3}{2};\dfrac{9}{2}\right)$}
	\loigiai{ 
		 Gọi $M(a;b;c)$.\\
		 Ta có $\overrightarrow{MA}=(1-a;1-b;1-c)$, $\overrightarrow{MB}=(5-a;-1-b;2-c)$, $\overrightarrow{MC}=(3-a;2-b;-4-c)$.\\
		$\overrightarrow{MA}+2\overrightarrow{MB}-\overrightarrow{MC}=\overrightarrow{0}\Leftrightarrow\heva{&1-a+2(5-a)-(3-a)=0\\&1-b+2(-1-b)-(2-b)=0\\&1-c+2(2-c)-(-4-c)=0}\Leftrightarrow \heva{&a=4\\&b=-\dfrac{3}{2}\\&c=\dfrac{9}{2}}\Rightarrow  M\left(4;-\dfrac{3}{2};\dfrac{9}{2}\right)$.
	} 
\end{ex}

% nguồn: Lop_12_Hinh_hoc_Chuong_2
\begin{ex}%[2H2H2-2]%[ma:b81b7a9e]
	Trong không gian $Oxyz$, hình  chiếu vuông góc của điểm $M(1; 2;-4)$ trên trục $Ox$~là điểm có tọa~độ
	\choice{$(0; 2;  0)$}
		{\True $(1;  0; 0)$}
		{$(0; 0;-4)$} 
		{$(1; 2;-4)$} 
	\loigiai{ 
		Hình chiếu vuông góc của điểm $M(1; 2;-4)$ trên trục  $Ox$~là $M'(1;0;0)$.
	 }
\end{ex}

% nguồn: Lop_12_Hinh_hoc_Chuong_2
\begin{ex}%[2H2H2-2]%[ma:393cd054]
	Trong không gian $Oxyz$, cho ba điểm $A(-4;1;-5)$, $B(2;-4;7)$, $C(3;-2;9)$. Tọa độ điểm $D$  để $ABCD$ là hình bình hành là
	 \choice
		 {$D(2;3;-3)$}
		{\True $D(-3;3;-3)$} 
		{$D(-3;-3;3)$} 
		{$D(-6;5;-12)$} 
	 \loigiai{
		Ta có  $\overrightarrow{BC}=(1;2;2)$.Tứ giác $ABCD$ là hình bình hành khi và chỉ khi
		\[\overrightarrow{AD}=\overrightarrow{BC}\Leftrightarrow\heva{&x_D+4=1\\&y_D-1=2\\&z_D+5=2}\Leftrightarrow\heva{&x_D=-3\\&y_D=3\\&z_D=-3.}\] 
	 }
\end{ex}

% nguồn: Lop_12_Hinh_hoc_Chuong_2
\begin{ex}%[2H2H2-2]%[ma:4b9e87d7]
	Trong không gian $Oxyz$, cho hình hộp $ABCD.A'B'C'D'$ có $A( 1;0;1),\,B( 2;1;2),\,D(  1;-1;1),\,C'( 4;5;-5)$. Tìm tọa độ tâm $I$của hình hộp.
	 \choice
		 {$I( -\dfrac{3}{2};\dfrac{5}{2};-3)$}
		{$I(  \dfrac{3}{2};\dfrac{5}{2};-3)$}
		 {$I( -\dfrac{5}{2};\dfrac{5}{2};-2)$}
		 {\True $I( \dfrac{5}{2};\dfrac{5}{2};-2)$}
	\loigiai{ 
		\begin{center}
			\begin{tikzpicture}[scale=0.7,font=\footnotesize,line join=round,line cap=round,>=stealth]
					\coordinate (A)  at (0,0); \coordinate (B) at (-2,-1.5); \coordinate (C) at (3,-1.5); \coordinate (D) at ($(A)+(C)-(B)$);
					\coordinate (A') at  ($(A)+(1,4)$);
					\coordinate  (B') at ($(A')+(B)-(A)$);
					\coordinate (C') at  ($(B')+(C)-(B)$);
					\coordinate  (D') at ($(A')+(D)-(A)$);
					\coordinate (I)  at ($(C)!0.5!(A')$);
					 \draw (D')--(A')--(B')--(C')--(D')--(D)--(C)--(B)--(B')
						(C)--(C'); 
					 \draw[dashed] (D)--(A)--(B)
						 (C)--(A')--(A)--(C');
					 \foreach \x/\g in {A/180,B/-90,C/-90,D/0,A'/120,B'/120,C'/90,D'/0,I/0} \fill (\x) circle (1pt)
						($(\x)+(\g:4mm)$)  node{$\x$};
				\end{tikzpicture}
		\end{center}
		Gọi $I$ là tâm  của hình hộp.\\
		Khi đó $I$ là trung  điểm của $AC'$, suy ra $I\left(\dfrac{5}{2};\dfrac{5}{2};-2\right)$.
	 }
\end{ex}

% nguồn: Lop_12_Hinh_hoc_Chuong_2
\begin{ex}%[2H2H2-2]%[ma:dfb31fae]
	Trong không gian $Oxyz$, cho ba điểm $A(1;2;-1)$, $B(2;-1;3)$, $C(-2;3;3)$. Điểm $M(a;b;c)$ là đỉnh thứ tư của hình bình hành $ABCM$.  Khi đó $P= a^2+b^2-c^2$ có giá trị bằng
	\choice 
		 {$42$}
		 {$-50$}
		 {$-48$}
		{\True $44$} 
	\loigiai{ 
		Ta có $\overrightarrow{BA} = (1-2, 2-(-1), -1-3) = (-1,3,-4)$.\\ 
		Ta có $\overrightarrow{CM} = (a-(-2), b-3, c-3) =  (a+2,b-3,c-3)$.\\
		Tứ giác $ABCM$ là hình bình hành  khi và chỉ khi $$\overrightarrow{BA} = \overrightarrow{CM}\Leftrightarrow\heva{&a+2=-1\\&b-3=3\\&c-3=-4} \Leftrightarrow \heva{&a=-3\\&b=6\\&c=-1.}$$
		Vậy  $M(-3;6;-1)$.\\
		Khi  đó $P=a^2+b^2-c^2=(-3)^2 + 6^2-(-1)^2=9+36-1=44$.
	} 
\end{ex}

% nguồn: Lop_12_Hinh_hoc_Chuong_2
\begin{ex}%[2H2H2-2]%[ma:91929de8]
	Trong không gian với hệ tọa độ $Oxyz$, cho hai điểm $A(1;1;0)$, $B(2;-1;3)$. Tìm tọa độ điểm  $C$ trên trục $Oy$ để tam giác $ABC$ vuông tại $A$.
	\choice 
		 {$\left(0;0;\dfrac{1}{2}\right)$}
		{$(0;2;0)$} 
		 {$\left(\dfrac{1}{2};0;0\right)$}
		{\True $\left(0;\dfrac{1}{2};0\right)$} 
	 \loigiai{
		Gọi $C(0;y;0)$,  $y\in \mathbb{R}$.
		Ta có $\overrightarrow{A  B}=(1;-2;3);\overrightarrow{A C}=(-1;y-1;0)$.\\
		Để tam giác $A B C$ vuông tại $A$, ta cần $$\overrightarrow{A B} \cdot \overrightarrow{A C}=0 \Leftrightarrow-1-2(y-1)+0=0 \Leftrightarrow-2 y+1=0  \Leftrightarrow y=\dfrac{1}{2}.$$
		Vậy tọa độ của  điểm $C$ là $C\left(0;\dfrac{1}{2};0\right)$.
	 }
\end{ex}

% nguồn: Lop_12_Hinh_hoc_Chuong_2
\begin{ex}%[2H2H2-2]%[ma:402c2602]
	Trong không gian với hệ tọa độ $Oxyz$, cho hai điểm $A(1;2;-1)$, $B(2;1;2)$. Điểm $M$ trên trục $Ox$ có hoành độ dương và thỏa mãn $M A^2+M B^2=23$. Khi đó tọa độ  điểm $M$ là
	 \choice
		 {\True $M(4;0;0)$}
		{$M(3;0;0)$} 
		 {$M(2;0;0)$}
		 {$M(1;0;0)$}
	\loigiai{Điểm $M$ thuộc trục $Ox$ có hoành độ dương suy ra  $M(a;0;0)$, $a>0$.\\
		 Ta có $\overrightarrow{M A}=(1-a;2;-1)$, $\overrightarrow{M B}=(2-a;1;2)$.\\
		Giả  thiết
		\allowdisplaybreaks 
		\begin{eqnarray*}
			M A^2+M B^2=23 &\Leftrightarrow&(1-a)^2+4+1+(2-a)^2+1+4=23\\&\Leftrightarrow& a^2-3 a-4=0 \Leftrightarrow\hoac{&a=-1&\text{ (loại)} \\&  a=4&\text{ (nhận)}.}
		\end{eqnarray*}
		Vậy  $M(4;0;0)$.
	 }
\end{ex}

% nguồn: Lop_12_Hinh_hoc_Chuong_2
\begin{ex}%[2H2H2-2]%[ma:8caba3e2]
	Trong không gia $Oxyz$, cho véctơ $\overrightarrow{a}=(-3;2;1)$ và điểm $A(4;6;-3)$. Tọa độ điểm  $B$ thỏa mãn $\overrightarrow{AB}=\overrightarrow{a}$ là
	 \choice
		 {$(-1;-8;2)$}
		{$(7;4;-4)$} 
		{\True $(1;8;-2)$} 
		{$(-7;-4;4)$} 
	 \loigiai{
		 Gọi $B(x;y;z)$. Khi đó $\overrightarrow{AB}=(x-4;y-6;z+3)$.\\
		Khi  đó $\overrightarrow{AB}=\overrightarrow{a}\Leftrightarrow\heva{&x-4=-3\\&y-6=2\\&z+3=1}\Leftrightarrow\heva{&x=1\\&y=8\\&z=-2.}$
	} 
\end{ex}

% nguồn: Lop_12_Hinh_hoc_Chuong_2
\begin{ex}%[2H2H2-2]%[ma:1632f63b]
	 Trong không gian với hệ tọa độ $Oxyz$, cho ba điểm $A(0;1;-2)$, $B(1;2;1)$, $C(4;3;m)$. Tìm $m$ để $4$ điểm $O$, $A$, $B$, $C$ đồng phẳng.
	 \choice
		 {$m=-7$}
		 {$m=-14$}
		{$m=7$} 
		 {\True $m=14$}
	 \loigiai{
		Để $4$ điểm $O$, $A$, $B$, $C$  đồng phẳng $\Leftrightarrow \overrightarrow{OC}=a\overrightarrow{OA}+b\overrightarrow{OB}\Leftrightarrow(4;3;m)=a(0;1;-2)+b(1;2;1)$\\
		 $\Rightarrow\heva{&4=b\\&3=a+2b\\&m=-2a+b}\Leftrightarrow\heva{&b=4\\&a=-5\\&m=14.}$
	} 
\end{ex}

% nguồn: Lop_12_Hinh_hoc_Chuong_2
\begin{ex}%[2H2H2-2]%[ma:e22e54c6]
	Trong không gian với hệ trục tọa độ $Oxyz$, cho ba điểm  $A(1;2;-1)$, $B(2;-1;3)$ và $C(-3;5;1)$. Điểm $D$ là đỉnh thứ tư của hình bình hành $ABCD$. Hãy chọn khẳng định đúng trong các khẳng định sau?
	\choice 
		{\True $D(-4;8;-3)$} 
		 {$D(-2;8;-3)$}
		 {$D(-4;8;-5)$}
		{$D(-2;2;5)$} 
	\loigiai{ 
		Giả sử  $D(x;y;z)$.\\
		Tứ giác $ABCD$ là hình bình hành khi  và chỉ khi $\overrightarrow{AB}=\overrightarrow{DC}\Leftrightarrow\heva{&-3-x=1\\&5-y=-3\\&1-z=4}\Rightarrow\heva{&x=-4\\&y=8\\&z=-3.}$\\
		Vậy $D(-4;8;-3)$. 
	 }
\end{ex}

% nguồn: Lop_12_Hinh_hoc_Chuong_2
\begin{ex}%[2H2H2-2]%[ma:30c051e2]
	Trong không gian $Oxyz$, cho tam giác $MNP$ có $M(2;-3;4), N(1;2;3)$ và $P(3;-2;2)$. Trọng tâm của tam giác $MNP$ có toạ  độ~là
	\choice 
		{\True  $(2;-1;3)$}
		 {$(6;-3;9)$}
		 {$(-2;1;-3)$}
		{$(-6;3;-9)$} 
	 \loigiai{
		Gọi $G(x;y;z)$~là trọng tâm của  tam giác $MNP$.\\
		 Suy ra
		 $\heva{&x=\dfrac{2+1+3}{3}=2\\&y=\dfrac{-3+2-2}{3}=-1\\&x=\dfrac{4+3+2}{3}=3.}$\\
		 Vậy $G(2;-1;3)$.
	 }
\end{ex}

% nguồn: Lop_12_Hinh_hoc_Chuong_2
\begin{ex}%[2H2H2-2]%[ma:7a27b9b4]
	Trong không gian với  hệ trục tọa độ $Oxyz$, điểm thuộc $Ox$ và cách đều hai điểm $A(4;2;-1)$ và $B(2;1;0)$ là
	\choice 
		 {$M(-4;0;0)$}
		{$M(5;0;0)$} 
		{\True  $M(4;0;0)$}
		{$M(-5;0;0)$} 
	\loigiai{Gọi $M(x;0;0)$ là điểm thuộc $Ox$ (với $x \in \mathbb{R}$).\\ 
		Điểm $M$ cách đều hai điểm $A(4;2;-1)$  và $B(2;1;0)$ khi và chỉ khi
		 \allowdisplaybreaks
		\begin{eqnarray*}
			M A=M B &\Leftrightarrow &\sqrt{(4-x)^2+2^2+(-1)^2}=\sqrt{(2-x)^2+1^2} \\&\Leftrightarrow & 21-8 x+x^2=5-4  x+x^2 \\&\Leftrightarrow & x=4.
		\end{eqnarray*}
		 Vậy $M(4;0;0)$.
	} 
\end{ex}

% nguồn: Lop_12_Hinh_hoc_Chuong_2
\begin{ex}%[2H2H2-2]%[ma:eb3c563b]
	Trong không gian $Oxyz$, hình  chiếu vuông góc của điểm $M(1; 2;-4)$ trên mặt phẳng $(Oxy)$~là điểm có tọa~độ
	\choice{$(1; 2;-4)$} 
		 {$(0; 2;-4)$}
		{$(1; 0;-4)$} 
		{\True $(1; 2;  0)$}
	\loigiai{ 
		Hình chiếu vuông góc của điểm $M(1; 2;-4)$ trên mặt  phẳng $(Oxy)$~là $M'(1;2;0)$.
	} 
\end{ex}

% nguồn: Lop_12_Hinh_hoc_Chuong_2
\begin{ex}%[2H2H2-2]%[ma:5c796b28]
	Cho tứ diện $SABC$ có $ABC$ là tam giác vuông tại $B$, $BC=3$, $BA=2$, $SA$ vuông góc với mặt phẳng $\left(ABC\right)$ và có độ dài bằng $2$. Với hệ toạ độ $Oxyz$ được thiết lập như hình bên (gốc tọa độ $O$ trùng với điểm $B$). Tìm khẳng  định sai
	\begin{center}
		\begin{tikzpicture}[scale=0.8,font=\footnotesize,line join=round,line cap=round,>=stealth]
				\coordinate (B)  at (0,0);
				\coordinate (A)  at (-2,2);
				\coordinate (C) at  (4,2);
				\coordinate (S) at  ($(A)+(0,3)$);
				\draw[dashed]  (A)--(C);
				\draw (A)--(S)--(C) 
					(S)--(B); 
				\foreach \i/\g in  {A/-170,B/-90,C/-90,S/90}{\draw[fill=blue](\i) circle (1.5pt)
					($(\i)+(\g:5mm)$)  node[scale=1]{$\i$};}
				 \coordinate (Z) at ($(B)+(0,5)$);
				 \coordinate (Y) at ($(B)!1.5!(A)$);
				\coordinate  (X) at ($(B)!1.5!(C)$);
				\draw[->] (B)--(Z); 
				\draw[->]  (B)--(X);
				\draw[->] (B)--(Y); 
				\draw (X)  node[scale=1,left]{$x$};
				 \draw (Y) node[scale=1,above]{$y$};
				 \draw (Z) node[scale=1,left]{$z$};
				\pic[draw,thin,angle radius=3mm] {right  angle = C--B--A};
				\pic[draw,thin,angle radius=3mm]  {right angle = C--B--Z};
				\pic[draw,thin,angle radius=2mm] {right angle  = Z--B--A};
			\end{tikzpicture}
	\end{center}
	 \choice
		 {$A\left(0; 2; 0\right)$}
		{$B\left(0;  0; 0\right)$}
		{$C\left(0; 0; 3\right)$} 
		 {\True $S\left(-2; 2; 2\right)$}
	\loigiai{ 
		Toạ độ  điểm $S\left(0; 2; 2\right)$.
	 }
\end{ex}

% nguồn: Lop_12_Hinh_hoc_Chuong_2
\begin{ex}%[2H2V2-2]%[ma:5627e64c]
	Trong không gian $Oxyz$, cho hai điểm $A(1;2;-1)$,  $B(2;-1;3)$. Tìm tọa độ điểm $I$ thỏa mãn $\overrightarrow{IA}+2\overrightarrow{IB}=\overrightarrow{0}$.
	 \choice
		{$I\left(-\dfrac{5}{3};0;\dfrac{5}{3}\right)$} 
		 {$I\left(\dfrac{5}{3};\dfrac{5}{3};0\right)$}
		{\True  $I\left(\dfrac{5}{3};0;\dfrac{5}{3}\right)$}
		 {$I\left(0;\dfrac{5}{3};\dfrac{5}{3}\right)$}
	\loigiai{ 
		Gọi $I(x;y;z)$.  Lúc dó $\overrightarrow{IA}=(1-x;2-y;-1-z)$, $\overrightarrow{IB}=(2-x;-1-y;3-z)$.\\
		 Suy ra $2\overrightarrow{IB}=(4-2x;-2-2y;6-2z)$ và $\overrightarrow{IA}+2\overrightarrow{IB}=(5-3x;-3y;5-3z)$.\\
		Từ đó có  $\overrightarrow{IA}+2\overrightarrow{IB}=\overrightarrow{0} \Leftrightarrow \heva{& 5-3x=0 \\& -3y=0 \\& 5-3z=0 } \Leftrightarrow \heva{& x=\dfrac{5}{3} \\& y=0 \\& z=\dfrac{5}{3}.}$\\
		Vậy  $I\left(\dfrac{5}{3};0;\dfrac{5}{3}\right)$.
	 }
\end{ex}

% nguồn: Lop_12_Hinh_hoc_Chuong_2
\begin{ex}%[2H2V2-2]%[ma:c591a135]
	Trong không gian  $Oxyz$, cho ba điểm $A(1;2;-1)$, $B(2;-1;3)$, $C(-4;7;5)$. Tọa độ chân đường phân giác trong góc $B$ của tam giác $ABC$ là
	 \choice
		{$(-2;11;1)$} 
		 {\True $\left(-\dfrac{2}{3};\dfrac{11}{3};1\right)$}
		{$\left(\dfrac{2}{3};\dfrac{11}{3};1\right)$} 
		 {$\left(\dfrac{11}{3};-\dfrac{2}{3};1\right)$}
	 \loigiai{
		Ta có  $BA=\sqrt{26}$, $BC=2\sqrt{26}$.\\
		Gọi $D$ là chân  đường phân giác trong góc $B$ thì $\dfrac{DA}{DC} = \dfrac{BA}{BC}=\dfrac{1}{2} \Rightarrow DC=2DA$.\\
		Vì $D$ là chân đường  phân giác trong nên $$2\overrightarrow{DA}+\overrightarrow{DC}=\overrightarrow{0} \Rightarrow \heva{&x_D=\dfrac{2x_A+x_C}{3}=-\dfrac{2}{3}\\&y_D=\dfrac{2y_A+y_C}{3}=\dfrac{11}{3}\\&z_D=\dfrac{2z_A+z_C}{3}=1.}$$
		 Vậy $D\left(-\dfrac{2}{3};\dfrac{11}{3};1\right)$.
	 }
\end{ex}

% nguồn: Lop_12_Hinh_hoc_Chuong_2
\begin{ex}%[2H2V2-2]%[ma:195ca647]
	Trong mặt phẳng tọa độ $Oxyz$, cho $A(0;2;-1)$, $B(2;1;5)$, $C(1;6;2)$ và hai điểm $M$, $N$ thỏa mãn  $\overrightarrow{MN}=\overrightarrow{NA}+\overrightarrow{NB}+\overrightarrow{NC}$. Đường thẳng $MN$ luôn đi qua điểm $I$ có tọa độ
	 \choice
		 {$(2;3;5)$}
		 {$(2;0;3)$}
		 {$(2;-1;3)$}
		{\True $(1;3;2)$} 
	\loigiai{ 
		Ta  có
		\begin{eqnarray*}
			 \overrightarrow{MN}&=&\overrightarrow{NA}+\overrightarrow{NB}+\overrightarrow{NC}\\ \overrightarrow{MN}&=&\overrightarrow{MA}+\overrightarrow{MB}+\overrightarrow{MC}-3\overrightarrow{MN}\\
			4\overrightarrow{MN}&=&\overrightarrow{MA}+\overrightarrow{MB}+\overrightarrow{MC}. 
		\end{eqnarray*}
		Gọi  $G$ là trọng tâm $\triangle ABC$, khi đó $4\overrightarrow{MN}=3\overrightarrow{MG}\Rightarrow \overrightarrow{MN}=\dfrac{3}{4}\overrightarrow{MG}$, suy ra $\overrightarrow{MN}$, $\overrightarrow{MG}$ cùng phương.\\
		Suy ra $M$, $G$,  $N$ thẳng hàng với $G(1;3;2)$.\\
		Vậy điểm $I(1;3;2)$ là  điểm cần tìm.
	 }
\end{ex}

% nguồn: Lop_12_Hinh_hoc_Chuong_2
\begin{ex}%[2H2V2-2]%[ma:9694e288]
	Trong không  gian $Oxyz$, cho hai điểm $A(1;3;4)$ và $B(1;0;1)$. Điểm $M$ nằm trên trục $Oz$ và cách đều hai điểm $A$, $B$ có tọa độ là
	\choice 
		{\True  $(0;0;4)$}
		{$(2;0;0)$} 
		{$(0;0;2)$} 
		 {$(0;4;0)$}
	\loigiai{ 
		Do điểm $M$ nằm  trên trục $Oz$ nên tọa độ của $M$ có dạng $M(0;0;z)$.\\
		Điểm $M$ cách đều hai điểm $A$, $B$ nên $MA = MB$ hay $MA^2=MB^2$.  Ta có
		\begin{align*}
			MA^2&=(1-0)^2+(3-0)^2+(4-z)^2=1+9+(4-z)^2=10+16-8z+z^2=z^2-8z+26,\\ 
			MB^2&=(1-0)^2+(0-0)^2+(1-z)^2=1+0+(1-z)^2=1+1-2z+z^2=z^2-2z+2. 
		\end{align*}
		Lúc đó  $MA^2 = MB^2 \Leftrightarrow z^2-8z+26=z^2-2z+2\Leftrightarrow6z=24 \Leftrightarrow z=4$.\\
		Vậy  $M(0;0;4)$.
	 }
\end{ex}

% nguồn: Lop_12_Hinh_hoc_Chuong_2
\begin{ex}%[2H2V2-2]%[ma:a35cca10]
	Trong không gian với hệ toạ độ $Oxyz$, cho ba điểm $A(-2;3;1)$, $B(2;1;0)$, $C(-3;-1;1)$. Tìm tất  cả các điểm $D$ sao cho $ABCD$ là hình thang có đáy $AD$ và $S_{ABCD}=3S_{\triangle ABC}$.
	\choice 
		 {$D(8;7;-1)$}
		{$\hoac{&D(-8;-7;1)  \\&D(12;1;-3)}$}
		 {$\hoac{&D(8;7;-1) \\&D(-12;-1;3)}$}
		{\True $D(-12;-1;3)$} 
	\loigiai{ 
		Gọi $D\left(x_D;y_D;z_D\right)$.\\ Lúc  đó $\overrightarrow{BC}=(-5;-2;1)$, $\overrightarrow{AD}=(x_D+2;y_D-3;z_D-1)$.\\
		 Ta có
		\begin{eqnarray*}
			S_{ABCD}=\dfrac{1}{2}(AD+BC)\cdot \mathrm{d}(A,BC)  &\Leftrightarrow& S_{ABCD}=\dfrac{1}{2}(AD+BC)\cdot \dfrac{2S_{\triangle ABC}}{BC}\\
			&\Leftrightarrow& 3S_{\triangle ABC}=\dfrac{(AD +  BC)\cdot S_{\triangle ABC}}{BC}\\
			&\Leftrightarrow& 3BC=AD +  BC\\
			&\Leftrightarrow& AD=2BC. 
		\end{eqnarray*}
		Mà $ABCD$ là hình thang có  đáy $AD$ nên
		$$\overrightarrow{AD}  = 2\overrightarrow{BC}\Leftrightarrow \heva{& x_D+2=-10\\&y_D-3=-4\\&z_D-1=2}\Leftrightarrow \heva{&x_D=-12\\&y_D=-1\\&z_D=3.}$$
		Vậy  $D(-12;-1;3)$.
	} 
\end{ex}

% nguồn: Lop_12_Hinh_hoc_Chuong_2
\begin{ex}%[2H2V2-2]%[ma:694f7d5c]
	Trong không gian $Oxyz$, cho ba điểm $E(1;3;2)$, $F(0;-1;5)$, $K(2;4;-1)$ và tam giác $ABC$ thỏa  mãn $\overrightarrow{AE}+\overrightarrow{BF}+\overrightarrow{CK}=\overrightarrow{0}$. Tọa độ trọng tâm $G$ của tam giác $ABC$ là
	\choice 
		{\True $G(1;2;2)$} 
		 {$G(-1;-4;3)$}
		{$G(2;2;1)$} 
		{$G(1;1;-3)$} 
	 \loigiai{
		Với $G$ là trọng tâm của tam giác $ABC$ thì ta  có $\overrightarrow{GA}+\overrightarrow{GB}+\overrightarrow{GC}=\overrightarrow{0}$.\\
		Lúc đó 
		\begin{eqnarray*}
			 &&\overrightarrow{AE}+\overrightarrow{BF}+\overrightarrow{CK}=\overrightarrow{0}\\
			&\Leftrightarrow& \left( \overrightarrow{AG}+\overrightarrow{GE}\right)+\left(\overrightarrow{BG}+\overrightarrow{GF}\right)+\left(\overrightarrow{CG}+\overrightarrow{GK}\right)=\overrightarrow{0}\\ 
			&\Leftrightarrow&  \left(\overrightarrow{GE}+\overrightarrow{GF}+\overrightarrow{GK}\right)-\left(\overrightarrow{GA}+\overrightarrow{GB}+\overrightarrow{GC}\right)=\overrightarrow{0}\\
			&\Leftrightarrow& \overrightarrow{GE}+\overrightarrow{GF}+\overrightarrow{GK}-\overrightarrow{0}=\overrightarrow{0}\\ 
			&\Leftrightarrow&  \overrightarrow{GE}+\overrightarrow{GF}+\overrightarrow{GK}=\overrightarrow{0}.
		\end{eqnarray*}
		Đẳng thức này chứng tỏ $G$ cũng là trọng  tâm của tam giác $EFK$.\\
		Do đó,  tọa độ của $G$ là
		$\heva{&x_G=\dfrac{1+0+2}{3}=1 \\& y_G = \dfrac{3+(-1)+4}{3}  = 2 \\&z_G = \dfrac{2+5+(-1)}{3} = 2.}$\\
		Vậy  $G(1;2;2)$.
	 }
\end{ex}

% nguồn: Lop_12_Hinh_hoc_Chuong_2
\begin{ex}%[2H2V2-2]%[ma:15f46371]
	Trong không gian với hệ tọa độ $Oxyz$, cho tứ diện $ABCD$ với  $A(1;-4;2)$, $B(2;1;-3)$, $C(3;0;-2)$ và $D(2;-5;-1)$. Điểm $G$ thỏa mãn $\overrightarrow{GA}+\overrightarrow{GB}+\overrightarrow{GC}+\overrightarrow{GD}=\overrightarrow{0}$ có tọa độ là
	 \choice
		{$G(2;-1;-1)$} 
		 {\True $G(2;-2;-1)$}
		{$G(0;-1;-1)$} 
		{$G(6;-3;-3)$} 
	\loigiai{ 
		Ta có 
		\begin{eqnarray*}
			\overrightarrow{GA}+\overrightarrow{GB}+\overrightarrow{GC}+\overrightarrow{GD}=\overrightarrow{0}&\Leftrightarrow& \overrightarrow{GO}+\overrightarrow{OA}+\overrightarrow{GO}+\overrightarrow{OB}+\overrightarrow{GO}+\overrightarrow{OC}+\overrightarrow{GO}+\overrightarrow{OD}=\overrightarrow{0}\\ 
			&\Leftrightarrow& 4\overrightarrow{OG}=\overrightarrow{OA}+\overrightarrow{OB}+\overrightarrow{OC}+\overrightarrow{OD}\\ &\Leftrightarrow&  \overrightarrow{OG}=\dfrac{1}{4}\left(\overrightarrow{OA}+\overrightarrow{OB}+\overrightarrow{OC}+\overrightarrow{OD}\right).
		\end{eqnarray*}
		 Gọi $G\left(x_G;y_G;z_G\right)$. Lúc đó
		 $$\heva{&x_G=\dfrac{1}{4}\left(x_A+x_B+x_C+x_D\right)= 2\\&y_G=\dfrac{1}{4}\left(y_A+y_B+y_C+y_D\right)= -2\\&z_G=\dfrac{1}{4}\left(z_A+z_B+z_C+z_D\right)=-1.}$$
		 Vậy $G(2;-2;-1)$.
	 }
\end{ex}

% nguồn: Lop_12_Hinh_hoc_Chuong_2
\begin{ex}%[2H2V2-2]%[ma:2a1eb89a]
	Trong không gian với hệ tọa độ $Oxyz$, cho hình vuông $ABCD$ trong đó $A(1;2;0)$, $B(3;0;8)$ và $C(-3;-6;8)$. Hai điểm $M$, $N$ lần  lượt nằm trên hai cạnh $AB$, $BC$ thỏa mãn $AM=BN=\dfrac{1}{3}\cdot BC$. Gọi $I(a;b;c)$ là giao điểm của $AN$, $DM$. Tính $P=a+b+c$.
	 \choice
		{$\dfrac{17}{5}$} 
		{\True $\dfrac{21}{5}$} 
		 {$\dfrac{20}{5}$}
		 {5}
	 \loigiai{
		\begin{center}
			\begin{tikzpicture}[
						 scale=1.0,
						line  cap=round,
						line  join=round]
					\coordinate  (A) at (0,0);
					\coordinate (B) at  (3,0);
					\coordinate  (C) at (3,3);
					\coordinate (D)  at (0,3);
					\coordinate (M) at ($(A)!0.33!(B)$); 
					\coordinate  (N) at ($(B)!0.33!(C)$);
					\draw (A) -- (B) -- (C) -- (D) -- cycle; 
					\draw[name path=path_AN]  (A) -- (N);
					\draw[name path=path_DM]  (D) -- (M);
					\path[name intersections={of=path_AN  and path_DM,by={I}}];
					\foreach \point/\position in {A/below right,B/below left,C/above  left,D/above right,M/below,N/right,I/above right}
						 {
					\fill (\point) circle  (0.5pt);
					\node[\position=2pt] at (\point) {$\point$}; 
					 }
				\end{tikzpicture}
		\end{center}
		Ta có  $\overrightarrow{AN}\cdot\overrightarrow{DM} = \left(\overrightarrow{AB} + \dfrac{1}{3}\overrightarrow{BC}\right)\cdot \left(\overrightarrow{DA} + \dfrac{1}{3}\overrightarrow{AB}\right) = 0 \Rightarrow AN \perp DM$.\\
		Lại có $\overrightarrow{BN} =  \dfrac{1}{3}\overrightarrow{BC}$, suy ra $N(1;-2;8)$.\\
		Vì hai tam giác $\triangle AIM$ và $\triangle ABN$ đồng dạng  nên
		$AI = \dfrac{AM\cdot AB}{AN}=\dfrac{1}{\sqrt{10}}AB$ và $AN = \dfrac{\sqrt{10}}{3}AB$ nên $AI =  \dfrac{3}{10}AN$.\\
		Suy ra $\overrightarrow{AI} = \dfrac{3}{10}\overrightarrow{AN} \Rightarrow I\left(1;\dfrac{4}{5};\dfrac{12}{5}\right) \Rightarrow P=1+\dfrac{4}{5}+\dfrac{12}{5}=\dfrac{5+4+12}{5}=\dfrac{21}{5}$. 
	 }
\end{ex}

% nguồn: Lop_12_Hinh_hoc_Chuong_2
\begin{ex}%[2H2V2-2]%[ma:c5f05ca3]
	Trong không gian  $Oxyz$, cho hai điểm $A(-4;1;5)$, $B(1;5;-3)$. Gọi $C$ là giao điểm của đường thẳng $AB$ và mặt phẳng $(Oyz)$. Trong các khẳng định sau, khẳng định nào đúng?
	 \choice
		{$\overrightarrow{AC}  = -\dfrac{1}{4}\overrightarrow{AB}$}
		 {$\overrightarrow{AB} = 5\overrightarrow{BC}$}
		 {$\overrightarrow{AC} = 4\overrightarrow{BC}$}
		{\True $\overrightarrow{AC} =  -4\overrightarrow{BC}$}
	\loigiai{ 
		Gọi  $C(0;a;b)$ là giao điểm của đường thẳng $AB$ và mặt phẳng $(Oyz)$.\\
		Ta có $\overrightarrow{AB}=(5;4;-8)$,  $\overrightarrow{AC}=(4;a-1;b-5)$.\\
		Vì $A$, $B$, $C$ thẳng hàng nên $\overrightarrow{AC}$ và $\overrightarrow{AB}$ cùng phương.\\ 
		Lúc  đó
		$$\dfrac{4}{5}=\dfrac{a-1}{4}=\dfrac{b-5}{-8} 
		\Leftrightarrow\heva{& a =  \dfrac{21}{5} \\& b=-\dfrac{7}{5}.}$$
		Từ  đó có $\overrightarrow{AC}=\left(4;\dfrac{16}{5};-\dfrac{32}{5}\right)$, $\overrightarrow{BC}=\left(-1;-\dfrac{4}{5};\dfrac{8}{5}\right)$.\\
		 Vậy $\overrightarrow{AC}=-4\overrightarrow{BC}$.
	 }
\end{ex}

% nguồn: Lop_12_Hinh_hoc_Chuong_2
\begin{ex}%[2H2V2-2]%[ma:86b2b86f]
	Trong không gian $Oxyz$, cho bốn điểm $A(2;0;0)$, $B(0;2;0)$, $C(0;0;2)$ và $D(2;2;2)$. Gọi $M$ và $N$ lần lượt là trung điểm của $AB$ và  $CD$. Tọa độ trung điểm $I$ của đoạn thẳng $MN$ là
	\choice 
		{$I(1;1;0)$} 
		{$I\left(\dfrac{1}{2};\dfrac{1}{2};1\right)$} 
		 {\True $I(1;1;1)$}
		 {$I(1;-1;2)$}
	 \loigiai{
		Vì $I(x;y;z)$  là trung điểm $MN$ nên ta có
		$$2\overrightarrow{IM} + 2\overrightarrow{IN} = \overrightarrow{0} \Leftrightarrow \overrightarrow{IA} + \overrightarrow{IB} + \overrightarrow{IC} +  \overrightarrow{ID} = \overrightarrow{0}\Leftrightarrow\overrightarrow{OI}=\dfrac{1}{4}\left(\overrightarrow{OA}+\overrightarrow{OB}+\overrightarrow{OC}+\overrightarrow{OD}\right).$$
		Lúc  đó
		 $$
		\heva{& 
		 x = \dfrac{2+0+0+2}{4} \\&
		y =  \dfrac{0+2+0+2}{4} \\&
		z =  \dfrac{0+0+2+2}{4}
		 }
		 \Leftrightarrow
		 \heva{&x=1\\&y=1\\&z=1.}
		$$ 
		 Vậy $I(1;1;1)$.
	 }
\end{ex}

% nguồn: Lop_12_Hinh_hoc_Chuong_2
\begin{ex}%[2H2V2-2]%[ma:ca1d35c6]
	Trong không gian $Oxyz$, cho hai điểm $A(4;5;6)$, $B(1;3;2)$. Gọi $H$ là hình chiếu vuông góc của $A$ lên mặt phẳng $(Oyz)$. Gọi $C$ là điểm nằm trên trục $Oz$ sao cho $BC$ và $AH$ là hai đường thẳng cắt nhau. Xác định tọa độ  điểm $C$.
	\choice 
		 {$C(0;0;2)$}
		 {$C\left(0;0;-\dfrac{2}{3}\right)$}
		{$C(0;0;4)$} 
		{\True $C(0;0;-4)$} 
	\loigiai{ 
		Vì $H$ là hình chiếu vuông  góc của $A$ lên mặt phẳng $(Oyz)$ nên $H(0;5;6)$.\\
		Gọi $C(0;0;c)$ là một điểm  nằm trên trục $Oz$ và $D$ là giao điểm của $BC$ và $AH$.\\
		Khi đó $A$, $D$, $H$  thẳng hàng nên $\overrightarrow{AD}=k\cdot \overrightarrow{HA}=(4k;0;0)$ suy ra tọa độ điểm $D(4k+4;5;6)$.\\
		Lại có  $B$, $D$, $C$ thẳng hàng nên
		$$\overrightarrow{BD}=l\cdot \overrightarrow{BC}\Leftrightarrow \heva{& 4k+4-1=l\cdot (0-1) \\& 5-3=l\cdot (0-3)  \\& 6-2=l\cdot (c-2)} \Leftrightarrow \heva{&k=-\dfrac{7}{12}\\&l=-\dfrac{2}{3}\\&c=-4 .}$$
		 Vậy $C(0;0;-4)$.
	 }
\end{ex}

% nguồn: Lop_12_Hinh_hoc_Chuong_2
\begin{ex}%[2H2V2-2]%[ma:13781bd9]
	Trong hệ trục tọa độ $Oxyz$, cho hai vectơ $\overrightarrow{PQ}=(0;1;-2)$, $\overrightarrow{PR}=(-2;-1;0)$ và điểm $M(1;-2;2)$ trung điểm của đoạn $QR$. Tọa độ điểm  $Q$ là
	 \choice
		{$(-1;1;-2)$} 
		{$(-2;2;-3)$} 
		 {$(0;1;3)$}
		 {\True $(2;-1;1)$}
	\loigiai{ 
		Ta có $\overrightarrow{RQ}=\overrightarrow{PQ}-\overrightarrow{PR}=(2;2;-2)$. Suy ra 
		$$\heva{&x_Q-x_R=2\\&y_Q-y_R=2\\&z_Q-z_R=-2.}\quad  (1)$$
		Mặt khác, điểm  $M(1;-2;2)$ trung điểm của đoạn $QR$ nên
		$$\heva{&x_Q+x_R=2\\&y_Q+y_R=-4\\&z_Q+z_R=4.}\quad (2)$$ 
		Từ $(1)$ và  $(2)$ suy ra $Q(2;-1;1)$.
	} 
\end{ex}


%%%%%%%%%%%%%%%% ĐÚNG/SAI %%%%%%%%%%%%%%%%

% nguồn: Lop_12_Hinh_hoc_Chuong_2
\begin{ex}%[2H2H2-2]%[ma:d0428697]
	\immini{ 
	Cho hình chóp $S.ABCD$ có đáy $ABCD$ là hình vuông có các cạnh bằng $1$, $SAD$ là tam giác đều và nằm trong mặt phẳng vuông góc với đáy. Gọi $O$, $M$ và $N$ lần lượt là  trung điểm của $AD$, $BC$ và $CD$. Thiết lập hệ trục tọa độ $Oxyz$ như hình vẽ.
	 \choiceTF
		{\True Tọa độ các  điểm $A$, $B$ là $A\left(0;-\dfrac{1}{2};0\right)$, $B\left(1;-\dfrac{1}{2};0\right)$}
		{\True Tọa độ các điểm  $C$, $D$ là $C\left(1;\dfrac{1}{2};0\right)$, $D\left(0;\dfrac{1}{2};0\right)$}
		{\True Tọa độ điểm  $S$ là $S\left(0;0;\dfrac{\sqrt{3}}{2}\right)$}
		{\True Tọa độ các điểm $M$, $N$  là $M(1;0;0)$, $N\left(\dfrac{1}{2};\dfrac{1}{2};0\right)$}
		}{\begin{tikzpicture}[scale=0.65,font=\footnotesize,line join=round,line cap=round,>=stealth]
		 \coordinate (A) at (0,0);
		\coordinate (D) at  (-2,-2);
		\coordinate (B) at  (5,0);
		\coordinate (C) at  ($(B)+(D)-(A)$);
		\coordinate (O) at ($(A)!0.5!(D)$); 
		\coordinate  (M) at ($(C)!0.5!(B)$);
		\coordinate  (S) at ($(O)+(0,5)$);
		 \coordinate (z) at ($(S)+(0,1)$);
		\coordinate (x)  at ($(M)+(1,0)$);
		\coordinate (y)  at ($(O)!2!(D)$);
		 \draw [->](S)--(z);
		 \draw [->](M)--(x);
		\draw [->](D)--(y); 
		\draw  (B)--(C)--(D);
		\draw  (S)--(D)
			(S)--(C) 
			 (S)--(B);
		 \draw[dashed](S)--(A)
			(A)--(D) 
			 (A)--(B)
			(S)--(O)-- (M); 
		 \pic[draw,angle radius=2mm] {right angle = S--O--B};
		\foreach  \i/\g in {S/135,A/45,D/-90,C/-90,B/45,O/-90,M/135}{\draw (\i) circle (1pt)
			($(\i)+(\g:4mm)$)  node[scale=1]{$\i$};}
		 \foreach \i/\g in {x/90,y/90,z/0}{\draw ($(\i)+(\g:3mm)$) node{$\i$};}
	\end{tikzpicture}
		} 
	\loigiai{ 
		Chọn  hệ tọa độ $Oxyz$ như hình vẽ.
		\begin{center}
			\begin{tikzpicture}[scale=0.8,font=\footnotesize,line join=round,line cap=round,>=stealth]
					 \coordinate (A) at (0,0);
					 \coordinate (D) at (-2,-2);
					\coordinate (B) at (5,0); 
					 \coordinate (C) at ($(B)+(D)-(A)$);
					 \coordinate (O) at ($(A)!0.5!(D)$);
					\coordinate  (M) at ($(C)!0.5!(B)$);
					\coordinate  (S) at ($(O)+(0,5)$);
					\coordinate (z) at ($(S)+(0,1)$); 
					\coordinate (x)  at ($(M)+(1,0)$);
					 \coordinate (y) at ($(O)!2!(D)$);
					\draw  [->](S)--(z);
					 \draw [->](M)--(x);
					 \draw [->](D)--(y);
					 \draw (B)--(C)--(D);
					\draw  (S)--(D)
						(S)--(C) 
						 (S)--(B);
					\draw[dashed](S)--(A) 
						(A)--(D) 
						(A)--(B) 
						 (S)--(O)-- (M);
					\pic[draw,angle radius=2mm] {right angle =  S--O--B};
					 \foreach \i/\g in {S/135,A/45,D/-90,C/-90,B/45,O/-90,M/135}{\draw (\i) circle (1pt)
						($(\i)+(\g:4mm)$) node[scale=1]{$\i$};} 
					\foreach \i/\g  in {x/90,y/90,z/0}{\draw ($(\i)+(\g:3mm)$) node{$\i$};}
				\end{tikzpicture}
		\end{center}
		$SAD$ là tam giác đều có cạnh bằng $1$ nên  $SO=\dfrac{\sqrt{3}}{2}$.\\
		$A\left(0;-\dfrac{1}{2};0\right)$, $B\left(1;-\dfrac{1}{2};0\right)$, $C\left(1;\dfrac{1}{2};0\right)$,  $D\left(0;\dfrac{1}{2};0\right)$, $S\left(0;0;\dfrac{\sqrt{3}}{2}\right)$, $M(1;0;0)$, $N\left(\dfrac{1}{2};\dfrac{1}{2};0\right)$.
		\begin{itemchoice}
			\itemch Tọa độ các  điểm $A$, $B$ là $A\left(0;-\dfrac{1}{2};0\right)$, $B\left(1;-\dfrac{1}{2};0\right)$.
			\itemch Tọa độ các  điểm $C$, $D$ là $C\left(1;\dfrac{1}{2};0\right)$, $D\left(0;\dfrac{1}{2};0\right)$.
			\itemch Tọa độ điểm $S$ là  $S\left(0;0;\dfrac{\sqrt{3}}{2}\right)$.
			\itemch Tọa độ các điểm $M$, $N$ là $M(1;0;0)$,  $N\left(\dfrac{1}{2};\dfrac{1}{2};0\right)$.
		\end{itemchoice}
	 }
\end{ex}

% nguồn: Lop_12_Hinh_hoc_Chuong_2
\begin{ex}%[2H2H2-2]%[ma:117998a6]
	 \immini[thm]{
	Cho ba điểm  $A(3;3;-6)$, $B(1;3;2)$ và $C(-1;-3;1)$. Gọi $M$, $N$, $K$ lần lượt là trung điểm của $AB$, $BC$ và $CA$.
	\choiceTF 
		{\True Tọa độ $M(2;3;2)$} 
		{Với $G$ là trọng tâm tam giác $ABC$ thì  $GC=2\sqrt{5}$}
		{\True Trọng  tâm tam giác $MNK$ là $E(1;1;-1)$}
		{\True Với $D(-3;-3;9)$ thì tứ giác $ABDC$ là hình bình hành} 
		 }
		 {
		\begin{tikzpicture}[scale=1,font=\footnotesize,line join=round,
					line cap=round,>=stealth] 
				\draw (0,0) coordinate (A)--(3,2.5) coordinate (B)--(4,0) coordinate (C)--cycle ($(A)!.5!(B)$)  coordinate (M)--($(C)!.5!(B)$) coordinate (N)--($(A)!.5!(C)$) coordinate (K)--cycle;
				\foreach  \x/\g in {A/180,B/90,C/0,M/140,N/45,K/-90}
				\fill  (\x) circle (1pt)
					($(\x)+(\g:3mm)$) node{$\x$}; 
			\end{tikzpicture}
		} 
	 \loigiai{
		\begin{itemchoice}
			\itemch $M$ là trung điểm của $AB$,  suy ra $M\left(\dfrac{x_A+x_B}{2};\dfrac{y_A+y_B}{2};\dfrac{z_A+z_B}{2}\right)$ hay $M(2;3;-2)$.
			\itemch Ta có $G(1;1;-1)$.  Suy ra $GC=\sqrt{(-1-1)^2+(-3-1)^2+(1+1)^2}=2\sqrt{6}$.
			\itemch Hai tam giác $ABC$ và $MNK$  có cùng trọng tâm. Suy ra $E$ trùng với $G(1;1;-1)$.
			\itemch  Ta có $\overrightarrow{AC}=(-4;-6;7)$, $\overrightarrow{BD}=(-4;-6;7)$ suy ra $\overrightarrow{AC}=\overrightarrow{BD}$.\\
			Vậy $ABDC$ là hình bình hành. 
		\end{itemchoice}
	 }
\end{ex}

% nguồn: Lop_12_Hinh_hoc_Chuong_2
\begin{ex}%[2H2H2-2]%[ma:c06c43a6]
	 Trong không gian $Oxyz$, cho hình hộp $ABCD.A'B'C'D'$ có $A(1; 0; 1)$, $B(2; 1; 2)$, $D(1;-1; 1)$, $D'(3;4;-6)$. Khi đó
	\choiceTF 
		{\True  Tọa độ $\overrightarrow{AB}=(1; 1; 1)$}
		{Tọa độ $C(2;  1; 2)$}
		{\True Tọa  độ $A'(3;5;-6)$}
		{Toạ  độ trọng tâm tam giác $A'B'C$ là $G(3; 4;-3)$}
	\loigiai{\begin{center}
		\begin{tikzpicture}[scale=0.6,font=\footnotesize,line join=round,line cap=round,>=stealth]
				 \path
					 (0,0) coordinate (A)
					++(-130:3)  coordinate (B)
					 ++(0:4) coordinate (C)
					($(A)+(C)-(B)$) coordinate  (D)
					 ($(A)!1/2!(C)$) coordinate (O);
				\foreach \i  in {A,B,C,D}{
					\coordinate (\i') at ($(\i)+(1,4)$); 
				 }
				 \draw (A')--(B')--(C')--(D')--cycle;
				 \draw (B)--(B')
					(C)--(C') 
					 (D)--(D')
					 (B)--(C)--(D);
				\draw[dashed,thin](B)--(A)--(A') 
					(A)--(D); 
				\foreach \i/\g in {A'/90,B'/90,C'/90,D'/90,A/-90,B/-90,C/-90,D/-90} 
				\fill  (\i) circle (1pt)+(\g:5mm) node[scale=1]{$\i$};
			\end{tikzpicture}
	\end{center}
	\begin{itemchoice}
		\itemch  Đúng.\\
		 Ta có $\overrightarrow{AB}=(1; 1; 1)$.
		\itemch Sai.\\ 
		Gọi $C(x; y; z)\Rightarrow \overrightarrow{DC}=(x-1;  y+1; z-1)$.\\
		Vì $ABCD$ là hình  bình hành\\
		$\Leftrightarrow \overrightarrow{DC}=\overrightarrow{AB}$ $\Leftrightarrow \heva{&{x-1=1} \\&{y+1=1} \\&{z-1=1}} \Leftrightarrow \heva{&{x=2} \\&{y=0} \\ &{z=2}}$  $\Rightarrow C(2;0;2)$.
		 \itemch Đúng.\\
		Ta có  $\overrightarrow{AD}=(0;-1; 0)$.\\
		Gọi $A'(x'; y'; z')\Rightarrow \overrightarrow{A'D'}=(3-x'; 4-y';-6-z')$.\\ 
		$ADD'A'$ là hình  bình hành\\
		$\Leftrightarrow \overrightarrow{A'D'}=\overrightarrow{AD}$ $\Leftrightarrow \heva{&{3-x'=0} \\&{4-y'=-1} \\&{-6-z'=0}} \Leftrightarrow \heva{&{x'=3} \\ &{y'=5} \\&{z'=-6}}$ $\Rightarrow  A'(3;5;-6)$.
		\itemch Sai.\\ 
		Gọi $B'\left(x_0; y_0; z_0\right)\Rightarrow  \overrightarrow{A'B'}=\left(x_0-3; y_0-5; z_0+6\right)$.\\
		$ABB'A'$ là  hình bình hành\\
		$\Leftrightarrow  \overrightarrow{A'B'}=\overrightarrow{AB}$ $\Leftrightarrow \heva{&{x_0-3=1} \\&{y_0-5=1} \\&{z_0+6=1}} \Leftrightarrow \heva{&{x_0=4} \\&{y_0=6} \\&{z_0=-5}} \Rightarrow B'(4; 6;-5)$.\\
		$G$ là trọng tâm tam giác $A'B'C$ $\Rightarrow \heva{&{x_G=\dfrac{2+3+4}{3}=3} \\&{y_G=\dfrac{0+5+6}{3}=\dfrac{11}{3}} \\&{z_G=\dfrac{2-6-5}{3}=-3}} \Rightarrow G\left(3;  \dfrac{11}{3};-3\right)$.
	\end{itemchoice}
 }
\end{ex}

% nguồn: Lop_12_Hinh_hoc_Chuong_2
\begin{ex}%[2H2H2-2]%[ma:c3a26ad9]
	Trong không gian,  cho ba điểm $A(4;2;1)$, $B(2;1;3)$, $C(-1;3;-2)$. Khi đó
	\choiceTF 
		{\True Tọa độ trọng tâm tam giác $ABC$  bằng $\left( \dfrac{5}{3};2;\dfrac{2}{3}\right)$}
		{\True Tọa độ trung  điểm đoạn thẳng $AB$ bằng $\left(3;\dfrac{3}{2};2\right)$}
		{\True Tứ giác $ABCD$ là  hình bình hành thì tọa độ điểm $D \left(1;4;-4\right)$}
		{Ba điểm $A$, $B$,  $C$ thẳng hàng}
	\loigiai{ 
		\begin{itemchoice}
			\itemch Tọa độ trọng tâm tam giác  $ABC$ là $\left( \dfrac{4+2-1}{3};\dfrac{2+1+3}{3};\dfrac{1+3-2}{3} \right)=\left( \dfrac{5}{3};2;\dfrac{2}{3}\right)$.
			\itemch Tọa độ trung điểm đoạn thẳng $AB$ là $\left(  \dfrac{4+2}{2};\dfrac{2+1}{2};\dfrac{1+3}{2}\right)=\left(3;\dfrac{3}{2};2\right)$.
			\itemch Để tứ giác $ABCD$ là hình bình hành ta cần 
			 $$
			\overrightarrow{AB}=\overrightarrow{DC}  \Leftrightarrow\heva{&2-4=-1-x_D\\&1-2=3-y_D\\&3-1=-2-z_D} \Leftrightarrow \heva{&x_D=1\\&y_D=4\\&z_D=-4}\Rightarrow D(1;4;-4).$$
			\itemch Ta có  $\overrightarrow{AB}=(-2;-1;2)$, $\overrightarrow{AC}=(-5;1;-3)$.\\
			Lúc đó, vì $\dfrac{-2}{-5} \neq \dfrac{-1}{1} \neq \dfrac{2}{-3}$ nên $\overrightarrow{AB}$ không cùng  phương với $\overrightarrow{AC}$.\\
			Suy ra  $A$, $B$, $C$ không thẳng hàng.
		\end{itemchoice}
	 }
\end{ex}

% nguồn: Lop_12_Hinh_hoc_Chuong_2
\begin{ex}%[2H2H2-2]%[ma:617e3900]
	Cho hình lăng  trụ $ABC\cdot A'B'C'$ có đáy $ABC$ là tam giác đều cạnh bằng $2$ như hình vẽ. Hình chiếu vuông góc của $A'$ lên $\left(ABC\right)$ trùng với trung điểm cạnh $AB$, góc $\widehat{A'AO}=60^{\circ}$. Với hệ toạ độ $Oxyz$ được thiết lập như hình bên (gốc tọa độ $O$ trùng với trung điểm của đoạn $BC$). hãy xét tính đúng sai của các khẳng định sau
	\begin{center}
		\begin{tikzpicture}[scale=0.8,font=\footnotesize,line join=round,line cap=round,>=stealth]
				 \coordinate (A) at (-2,0);
				\coordinate (B) at (-0.5,-2); 
				\coordinate (C)  at (2,0);
				\coordinate (H) at (1,4); 
				 \coordinate (O) at ($(A)!0.5!(B)$);
				\foreach \i in {A,B,C}{\coordinate (\i')  at ($(\i)+(H)$); \draw (\i)--(\i');}
				 \draw (A')--(B')--(C')--cycle;
				\draw  (A)--(B)--(C)
					(O)--(A'); 
				\draw[dashed,thin](A)--(C)--(O); 
				\pic[draw,thin,angle radius=2mm] {right angle = A'--O--A}; 
				\pic[draw,thin,angle  radius=2mm] {right angle = B--O--C};
				\foreach \i/\g in  {A/180,B/-90,C/0,A'/180,B'/-30,C'/0,O/-100}{\draw (\i) circle (1.5pt)
					($(\i)+(\g:4mm)$) node[scale=1]{$\i$};} 
				\coordinate  (Z) at ($(O)!1.2!(A')$);
				 \coordinate (Y) at ($(O)!1.5!(C)$);
				\coordinate (X) at  ($(O)!2!(B)$);
				\draw[->]  (A')--(Z);
				\draw[->] (C)--(Y); 
				 \draw[->] (B)--(X);
				\draw (X)  node[scale=1,left]{$x$};
				\draw  (Y) node[scale=1,above]{$y$};
				\draw  (Z) node[scale=1,left]{$z$};
				\pic[draw,thin,angle radius=3mm] {right angle  = C--O--A'};
			\end{tikzpicture}
	\end{center}
	 \choiceTF
		{\True Tọa độ điểm $A\left(-1;0;0\right)$} 
		{\True Tọa độ điểm $C\left(0;\sqrt{3};0\right)$} 
		{Tọa độ điểm  $A'\left(0;-1;\sqrt{3} \right)$}
		{\True Tọa độ  điểm $C'\left(1;\sqrt{3};\sqrt{3} \right)$}
	\loigiai{ 
		Độ dài $OC=2\cdot \dfrac{\sqrt{3}}{2}=\sqrt{3} \cdot OA'=OA\cdot \tan60^{\circ}=\sqrt{3}$. Với  hệ trục $Oxyz$ đã chọn như hình vẽ trên thì
		\begin{itemchoice}
			\itemch Điểm  $A\in Ox$, nằm ngược chiều dương và $OA=1$ nên $A\left(-1;0;0\right)$.
			\itemch  Điểm $A'\in Oy$, nằm cùng chiều dương và $OC=\sqrt{3}$ nên $C\left(0;\sqrt{3};0\right)$.
			\itemch $A'\in Oz$, nằm  cùng chiều dương và $OA'=\sqrt{3}$ nên $A'\left(0;0;\sqrt{3} \right)$.
			\itemch Gọi $C'\left(x;y;z\right)$. Ta có $\overrightarrow{A'C'}=\overrightarrow{AC}\Leftrightarrow \heva{x-0&=1 \\ y-0&=\sqrt{3} \\ z-\sqrt{3}&=0}  \Leftrightarrow \heva{x&=1 \\ y&=\sqrt{3}\\ z&=\sqrt{3}}$.
		\end{itemchoice}
	 }
\end{ex}

% nguồn: Lop_12_Hinh_hoc_Chuong_2
\begin{ex}%[2H2H2-2]%[ma:d2269dc1]
	Trong không gian  $Oxyz$, cho ba điểm $A(3; 5; -1)$, $B(7; x; 1)$, $C(9; 2; y)$, $D(-1; 2; 3)$.
	\choiceTF 
		{Toạ độ điểm $M$ trung điểm  của đoạn $AD$ là $(1; 7; 1)$}
		{\True Toạ độ điểm $E$ để tứ  giác $OAD$E là hình bình hành là $(-4; -3; 4)$}
		{Toạ độ điểm $C$ để tam  giác $ACD$ vuông tại $A$ là $(9; 2; 11)$}
		{\True Để ba điểm $A$, $B$, $C$ thẳng hàng thì giá trị của biểu thức  $x+y$ bằng $5$}
	\loigiai{ 
		\begin{itemchoice}
			\itemch $M$  là trung điểm đoạn $AD$ nên\\ $M = \left( \dfrac{3 + (-1)}{2}; \dfrac{5 + 2}{2}; \dfrac{-1 + 3}{2} \right) = \left(1; \dfrac{7}{2}; 1\right) \neq (1; 7; 1)$.
			\itemch Giả sử $E(x; y; z)$. Để $OAD$E là hình bình hành \\$\Leftrightarrow \overrightarrow{OE} = \overrightarrow{AD} = (-4; -3; 4) \Rightarrow$ 
			 $\heva{
			 &x = -4 \\
			 &y = -3 \\
			&z =  4
			 }$.\\
			Vậy $E =  (-4; -3; 4)$.
			\itemch $\overrightarrow{AC} = (6;  -3; y+1)$, $\overrightarrow{AD} = (-4; -3; 4)$.\\
			Để tam giác $ACD$ vuông tại $A$ thì\\ 
			$AC \perp  AD$\\
			$\Leftrightarrow \overrightarrow{AC} \cdot \overrightarrow{AD} =  0$\\
			$\Leftrightarrow 6(-4)  + (-3)(-3) + 4(y+1) =0$ \\
			$\Leftrightarrow -24 + 9 + 4y + 4 = 4y  - 11 = 0$ \\
			$\Rightarrow y =  \dfrac{11}{4}$\\
			$\Rightarrow C = (9; 2; \dfrac{11}{4}) \neq (9;  2; 11)$.
			\itemch $\overrightarrow{AB} = (4; x - 5; 2)$, $\overrightarrow{AC} = (6; -3; y + 1)$. Nếu $A$, $B$, $C$ thẳng  hàng thì\\
			$\overrightarrow{AB}  = k \cdot \overrightarrow{AC} \Rightarrow$
			 $\heva{
			&4 = 6k \\ 
			&x - 5  = -3k \\
			&2 = k(y + 1) 
			 } \Rightarrow
			 \heva{
			&k = \dfrac{2}{3}  \\
			 &x = 3 \\
			&y =  2.
			} \Rightarrow x  + y = 5$.
		\end{itemchoice}
	} 
\end{ex}


%%%%%%%%%%%%%%%% TRẢ LỜI NGẮN %%%%%%%%%%%%%%%%

% nguồn: Lop_12_Hinh_hoc_Chuong_2
\begin{bt}%[2H2N2-2]%[ma:bf3bbf43]
	Trong không gian với hệ tọa~độ $Oxyz$, cho  hai điểm $A(6; 4;-1),$ $ B(0; 0; 3)$ và $M(a; b; c)$~là trung điểm của đoạn thẳng $AB$. Giá trị của biểu thức $a+b+c$~là bao nhiêu?
	 \shortans{6}
	\loigiai{Ta có $M$~là trung điểm $AB$ nên $$\heva{&x_M=\dfrac{x_A+x_B}{2}\\&y_M=\dfrac{y_A+y_B}{2}\\&z_M=\dfrac{z_A+z_B}{2}}\Leftrightarrow \heva{&a=\dfrac{6+0}{2}=3\\&b=\dfrac{4+0}{2}=2\\&c=\dfrac{-1+3}{2}=1}\Rightarrow  a+b+c=3+2+1=6.$$
	 }
\end{bt}

% nguồn: Lop_12_Hinh_hoc_Chuong_2
\begin{bt}%[2H2N2-2]%[ma:6787feae]
	Trong không gian với hệ tọa~độ $Oxyz$, cho tam giác $ABC$ có ba  đỉnh $A(2; 1;-3)$, $B(4; 2; 1),$ $ C(3; 0; 5)$ và $G(a; b; c)$~là trọng tâm của tam giác $ABC$. Giá trị của biểu thức $P= abc$~là bao nhiêu?
	 \shortans{3}
	\loigiai{Ta có $G$~là trọng  tâm tam giác $ABC$ nên $$\heva{&x_G=\dfrac{x_A+x_B+x_C}{3}\\&y_G=\dfrac{y_A+y_B+y_C}{3}\\&z_G=\dfrac{z_A+z_B+z_C}{3}}\Leftrightarrow \heva{&a=\dfrac{2+4+3}{3}=3\\&b=\dfrac{1+2+0}{3}=1\\&c=\dfrac{-3+1+5}{3}=1}\Rightarrow abc=3\cdot 1\cdot 1=3.$$
	} 
\end{bt}

% nguồn: Lop_12_Hinh_hoc_Chuong_2
\begin{bt}%[2H2H2-2]%[ma:b36b0956]
	 Cho ba điểm $A(2; 5; 3)$, $B(3; 7; 4)$, $C(a; b; 6)$. Biết ba điểm $A$, $B$, $C$ thẳng hàng. Tính $2a+b$.
	\shortans{21} 
	 \loigiai{
		Ta có  $\overrightarrow{AB}=(1; 2; 1)$, $\overrightarrow{AC}=(a-2; b-5; 3)$.\\
		Ba điểm $A$, $B$, $C$ thẳng  hàng khi và chỉ khi $\overrightarrow{AB}$, $\overrightarrow{AC}$ cùng phương.\\
		Suy ra $\dfrac{a-2}{1}=\dfrac{b-5}{2}=\dfrac{3}{1}$.\\ 
		 Ta có $a-2=3\Leftrightarrow a=5$ và
		$\dfrac{b-5}{2}=3\Leftrightarrow  b=11$.\\
		Suy ra $2a+b=2\cdot  5+11=21$.
	} 
\end{bt}

% nguồn: Lop_12_Hinh_hoc_Chuong_2
\begin{bt}%[2H2H2-2]%[ma:a6d6e392]
	Trong không gian với hệ tọa độ $Oxyz$, cho tam giác $ABC$ với $A(1;-3;3)$; $B(2;-4;5)$; $C(a;-2;b)$ nhận điểm $G(1;c;6)$ làm trọng  tâm của nó thì giá trị của tổng $a+b+c$ bằng
	 \shortans{7}
	\loigiai{ 
		Vì $G$ là trọng tâm tam giác $ABC$  nên ta có
		 $\heva{&x_A+x_B+x_C=3x_G\\&
		 y_A+y_B+y_C=3y_G\\&
		z_A+z_B+z_C=3z_G}$\\ 
		$\Rightarrow\heva{&1+2+a=3\\&-3-4-2=3c\\&3+5+b=18}\Rightarrow\heva{&a=0\\&b=10\\&c=-3}\Rightarrow  a+b+c=7$.
	 }
\end{bt}

% nguồn: Lop_12_Hinh_hoc_Chuong_2
\begin{bt}%[2H2H2-2]%[ma:520a62e3]
	Trong không gian  $Oxyz$, cho hai điểm $A(3;1;-2)$, $B(2;-3;5)$. Điểm $M(x;y;z)$ thuộc đoạn $AB$ sao cho $MA=2MB$. Tính $x+3y+z$
	\shortans{0} 
	\loigiai{ 
		Vì $M$ thuộc đoạn $AB$ nên $\overrightarrow{MA}=-2\overrightarrow{MB}\Leftrightarrow\heva{&3-x=-2(2-x)\\&1-y=-2(-3-y)\\&-2-z=-2(5-z)}\Leftrightarrow\heva{&x=\dfrac{7}{3}\\&y=-\dfrac{5}{3}\\&z=\dfrac{8}{3}}$.\\ 
		Vậy $x+3y+z=\dfrac{7}{3}+3\cdot  \left(-\dfrac{5}{3}\right)+ \dfrac{8}{3}=0$.
	 }
\end{bt}

% nguồn: Lop_12_Hinh_hoc_Chuong_2
\begin{bt}%[2H2H2-2]%[ma:5a7b8344]
	Trong không gian với hệ tọa độ $Oxyz$, cho tam  giác $ABC$ với $A\left(1;-3;3\right)$, $B\left(2;-4;5\right)$, $C\left(a;-2;b\right)$ nhận điểm $G\left(1;c;6\right)$ làm trọng tâm của nó thì giá trị của tổng $a+b+c$ bằng
	 \shortans{7}
	\loigiai{ 
		Vì $G$ là trọng  tâm tam giác $ABC$ nên ta có
		\begin{eqnarray*}
			&&\heva{&x_A+x_B+x_C=3x_G  \\&y_A+y_B+y_C=3y_G \\&z_A+z_B+z_C=3z_G}\\
			&\Rightarrow&\heva{&1+2+a=3  \\&-3-4-2=3c \\&3+5+b=18}\\
			&\Rightarrow&\heva{&a=0 \\&b=10 \\&c=-3.} 
		\end{eqnarray*}
		Do  đó $a+b+c=7$.
	 }
\end{bt}

% nguồn: Lop_12_Hinh_hoc_Chuong_2
\begin{bt}%[2H2H2-2]%[ma:319056f5]
	Trong không gian với hệ tọa độ $Oxyz$ cho hai điểm $A(4;2;1)$, $B(-2;-1;4)$. Tìm được tọa độ điểm $M(a;b;c)$  thỏa mãn đẳng thức $\overrightarrow{AM}=2\overrightarrow{MB}$. Tính $a+b+c$
	\shortans{3} 
	 \loigiai{
		Gọi điểm  $M(a;b;c)$. Khi đó $\overrightarrow{AM}=2\overrightarrow{MB}\Leftrightarrow\heva{&a-4=2(-2-a)\\&b-2=2(-1-b)\\&c-1=2(4-c)}\Leftrightarrow\heva{&a=0\\&b=0\\&c=3.}$\\
		Vậy  $a+b+c=3$.
	 }
\end{bt}

% nguồn: Lop_12_Hinh_hoc_Chuong_2
\begin{bt}%[2H2H2-2]%[ma:aa6c1c7d]
	 \immini[thm]{
	Để chuẩn bị cho ngày hội thao, người ta dựng bốn chiếc cột tại bốn góc của một sân bóng hình chữ nhật  với kích thước là $15\ \mathrm{m} \times 25\ \mathrm{m}$. Bốn chiếc cột vuông góc với mặt sân và có chiều cao lần lượt là $3$ mét, $4$ mét, $6$ mét và $c$ mét. Một tấm bạt lớn được căng phẳng với bốn góc được cố định vào đầu bốn cột. Xét hệ trục tọa độ $Oxyz$ như hình vẽ bên (đơn vị trên các trục là mét) thì điểm $D'$ có tọa độ là $(a; b; c)$. Tính $a-2 b+c$.
	 \shortans{-45}
	}{ 
	\begin{tikzpicture}[>=stealth,line join=round,line cap=round,font=\footnotesize,scale=1]
			 \path
				 (0,0) coordinate (A)
				(0:5)  coordinate (D)
				(90:0.6)  coordinate (A')
				(-120:3) coordinate (B) 
				($(B)+(D)-(A)$)  coordinate (C)
				($(B)+(90:0.8)$) coordinate (B') 
				 ($(C)+(90:1.2)$) coordinate (C')
				 ($(D)+(90:1.2)$) coordinate (D')
				(-120:4) coordinate (x) 
				(6,0)  coordinate (y)
				 (0,2) coordinate (z);
			 \draw[->] (D)--(y) node[below right]{$y$};
			\draw[->] (B)--(x) node[below right]{$x$}; 
			\draw[->] (A')--(z) node[below right]{$z$}; 
			 \draw
				(A')--(B')--(C')--(D')--cycle  (B')--(B) node[midway,left]{\scriptsize 4m}--(C) node[midway,below]{\scriptsize 25m}--(C') node[midway,left]{\scriptsize 6m}
				(C)--(D) node[midway,right]{\scriptsize 15m}--(D') node[midway,right]{\scriptsize $c$}; 
			\node at ($(A)+(-35:0.5)$){\scriptsize  $O \equiv A$};
			\draw[dashed]  (A')--(A) node[midway,right]{\scriptsize 3m}--(B)
				(A)--(D); 
			\foreach \x/\g in {B/180,C/-45,D/45,A'/180,B'/135,C'/0,D'/45}\fill (\x) circle  (1pt) +(\g:3mm) node {$\x$};
		\end{tikzpicture}
	 }
	\loigiai{ 
		Theo hệ trục tọa độ đã chọn, ta có $A'(0,0,3)$;  $B'(15,0,4)$; $C'(15,25,6)$.\\
		Vì $A'B'C'D'$ là  hình bình hành nên
		\[\heva{&x_{D'}=x_{A'}+x_{C'}-x_{B'}=0+15-15=0\\&y_{D'}=y_{A'}+y_{C'}-y_{B'}=0+25-0=25\\&z_{D'}=z_{A'}+z_{C'}-z_{B'}=3+6-4=5.}  \]
		Suy ra $a=0$, $b=25$ và $c=5$.\\ 
		Vậy $a-2b+c=0-2\cdot  25+5=45$.
	} 
\end{bt}

% nguồn: Lop_12_Hinh_hoc_Chuong_2
\begin{bt}%[2H2V2-2]%[ma:0e68aae0]
	Trong không gian $Oxyz$, cho ba điểm $A(-1; 1; 2)$, $B(3; 2;  2)$, $C(-1; 6; 0)$. Xét $M(a; b; 0)$ trên mặt phẳng $\left(Oxy\right)$ sao cho biểu thức $S= 2MA + \left| \overrightarrow{MB} +\overrightarrow{MC}\right|$ đạt giá trị nhỏ nhất. Khi đó $a+b$ bằng bao nhiêu (kết quả làm tròn đến hàng phần chục)?
	 \shortans{3,3}
	 \loigiai{
		\begin{center}
			\begin{tikzpicture}[scale=1.4,>=stealth,font=\footnotesize,line join=round,line cap=round]
					 \def\d{6}
					 \def\r{3}
					\path (0:0) coordinate (F) 
						++(0:\d) coordinate  (N)
						++(50:\r)  coordinate (K)
						($(F)+(K)-(N)$) coordinate  (E)
						(40: 2) coordinate  (A)
						 (44:3) coordinate (B)
						(20:5) coordinate (C) 
						($(B)!.5!(C)$) coordinate  (I)
						($(A)+(-90:2.5)$)  coordinate (A')
						($(A')!.7!(I)$)  coordinate (M)
						(intersection of A--A' and  F--N) coordinate (Td)
						($(A)+(-90:.6)$) coordinate (Mp); 
					 \draw[dashed] (Mp)--(Td);
					 \draw (F)--(N)--(K)--(E)--cycle
						 (A)--(Mp)
						(B)--(C) 
						 (I)--(A')--(Td)
						pic[draw,"$Oxy$",angle radius=1.2cm]{angle=N--F--E}; 
					\foreach \x/ \goc in {A/0,B/-40,C/0,I/50,A'/30,M/0} 
					\fill (\x)  circle (1pt)
						($(\x)+(\goc:3mm)$)  node {$\x$};
				\end{tikzpicture}
		\end{center}
		Trung điểm của $BC$  là $I(1;4;1)$.\\
		 Ta có $S=2MA+\left|2\overrightarrow{MI}\right|=2\left(MA + MI\right)$.\\
		Do điểm $A$ và điểm $I$ nằm cùng phía so với mặt phẳng $\left(Oxy\right)$ nên $MA + MI  = MA'+MI$ với $A'(-1;1;-2)$ là điểm đối xứng của $A$ qua mặt phẳng $\left(Oxy\right)$.\\
		Ta  lại có $S= 2\left(MA'+MI\right)\geq 2A'I$.\\
		Suy ra  $M \in \left(Oxy\right)$ và $M$, $A'$, $I$ thẳng hàng.\\
		Khi đó, hai vectơ $\overrightarrow{A'M}=(a+1; b-1; 2),  \overrightarrow{A'I}=(2; 3; 3)$ cùng phương.\\
		Suy ra $\dfrac{a+1}{2}=\dfrac{b-1}{3}=\dfrac{2}{3} \Rightarrow \heva{&a=\dfrac{1}{3}\\ & b=3.}$\\ 
		Vậy $a+b=\dfrac{10}{3} \approx  3{,}3$.
	} 
\end{bt}

% nguồn: Lop_12_Hinh_hoc_Chuong_2
\begin{bt}%[2H2V2-2]%[ma:f161c861]
	Trong không gian $Oxyz$, cho ba điểm $A(1;4;5)$, $B(3;4;0)$, $C(2;-1;0)$ và mặt phẳng $(P)\colon 3x+3y-2z-29=0$. Gọi $M(a;b;c)$ là  điểm thuộc $(P)$ sao cho biểu thức $T=MA^2+MB^2+3MC^2$ đạt giá trị nhỏ nhất. Tính tổng $a+b+c$.
	 \shortans{8}
	\loigiai{ 
		Chọn điểm $I\left(x_0; y_0; z_0\right)$ thỏa  mãn $\overrightarrow{IA}+\overrightarrow{IB}+3\overrightarrow{IC} = \overrightarrow{0}$. Ta có
		\begin{align*}
			 &\overrightarrow{IA}=\left(1-x_0; 4-y_0; 5-z_0\right),
			 \overrightarrow{IB}=\left(3-x_0; 4-y_0; -z_0\right),
			\overrightarrow{IC}=\left(2-x_0; -1-y_0;  -z_0\right). \\
			 \Rightarrow &\overrightarrow{IA}+\overrightarrow{IB}+3\overrightarrow{IC} = \overrightarrow{0}
			 \Leftrightarrow \heva{&10-5x_0=0 \\ &5-5y_0=0 \\ &5-5z_0=0}
			 \Leftrightarrow \heva{&x_0=2 \\ &y_0=1 \\ &z_0=1}
			\Rightarrow I(2;  1; 1).
		\end{align*}
		Ta có 
		\begin{align*}
			T &= MA^2 + MB^2 + 3MC^2  \\
			&= \overrightarrow{MA}^2 + \overrightarrow{MB}^2 +  3\overrightarrow{MC}^2 \\
			&= \left(\overrightarrow{MI}+\overrightarrow{IA}\right)^2 + \left(\overrightarrow{MI}+\overrightarrow{IB}\right)^2 +  3\left(\overrightarrow{MI}+\overrightarrow{IC}\right)^2 \\
			&= 5MI^2 + 2\overrightarrow{MI}.\left(\overrightarrow{IA}+\overrightarrow{IB}+3\overrightarrow{IC}\right) + IA^2 + IB^2  + 3IC^2 \\
			 &= 5MI^2 + IA^2 + IB^2 + 3IC^2.
		\end{align*}
		Từ đó ta có $T_{\min} \Leftrightarrow  MI_{\min}$. Suy ra $M$ là hình chiếu vuông góc của $I$ trên mặt phẳng $(P)$. \\
		Gọi $\Delta$ là đường thẳng đi qua $I$ có vectơ chỉ phương là $\overrightarrow{u}_{\Delta}$ và vuông góc với $(P)$ có vectơ pháp tuyến là $\overrightarrow{n}_{(P)}$ \\ 
		$\Rightarrow \overrightarrow{u}_{\Delta} =  \overrightarrow{n}_{P} = (3;3;-2)
		\Rightarrow \Delta\colon  \heva{&x=2+3t \\ &y=1+3t \\ &z=1-2t.}$ \\
		$M  \in \Delta \Rightarrow M(2+3t;1+3t;1-2t)$.\\
		$M \in (P) \Leftrightarrow 3(2+3t) + 3(1+3t) -  2(1-2t) - 29 = 0
		\Leftrightarrow  t=1 \Rightarrow M(5;4;-1)$.\\
		Vậy $a+b+c =  5+4-1=8$.
	 }
\end{bt}

% nguồn: Lop_12_Hinh_hoc_Chuong_2
\begin{bt}%[2H2V2-2]%[ma:f0d2bed7]
	Trong không gian $Oxyz$, cho tam  giác $ABC$ có điểm $A(1; 2; 3)$, $B(2;-1; 3)$ và $C(-1; 1; 1)$. Biết rằng toạ độ của chân đường cao hạ từ $A$ xuống $BC$ là $H\ (a; b; c)$, giá trị của $P=17(a+b+c)$ bằng
	 \shortans{41}
	\loigiai{ 
		Ta có $\overrightarrow{AH}=(a-1; b-2; c-3)$, $\overrightarrow{BC}=(-3; 2;-2)$, $\overrightarrow{BH}=(a-2;  b+1; c-3)$.\\
		Vì $H$ là chân  đường cao hạ từ $A$ xuống $BC$ nên ta có $AH\perp BC$ và $B$, $H$, $C$ thẳng hàng. Từ đây suy ra
		\allowdisplaybreaks 
		\begin{eqnarray*}
			\heva{&\overrightarrow{AH} \cdot \overrightarrow  {BC}=0 \\
			&\overrightarrow{BH}=k \overrightarrow {BC}} 
			 \Leftrightarrow \heva{&-3(a-1)+2(b-2)-2(c-3)=0 \\
			 &a-2=-3k \\
			&b+1=2k \\ 
			 &c-3=-2k}
			\Leftrightarrow  \heva{&a=\dfrac{7}{17} \\
			&b=\dfrac{1}{17} \\ 
			&c=\dfrac{33}{17}  \\
			&k=\dfrac{9}{17}. 
			} 
		\end{eqnarray*}
		Do đó  $P=17(a+b+c)=41$.
	 }
\end{bt}

% nguồn: Lop_12_Hinh_hoc_Chuong_2
\begin{bt}%[2H2V2-2]%[ma:a40c7311]
	Cho  hai điểm $A(3; -1; 2)$, $B(2; 3; -3)$, $C(-2; 1; -2)$ và mặt phẳng $(Oyz)$. Gọi $M(a; b; c)$ là điểm thuộc mặt phẳng $(Oyz)$ sao cho biểu thức
	 $\overrightarrow{MA} \cdot \overrightarrow{MB} + \overrightarrow{MB} \cdot \overrightarrow{MC} + \overrightarrow{MC} \cdot \overrightarrow{MA}$ đạt giá trị nhỏ nhất. Tính $a - 2b + c$.
	\shortans{-3} 
	 \loigiai{
		Gọi $I(a; b; c)$ sao cho $\overrightarrow{IA} + \overrightarrow{IB} + \overrightarrow{IC} = \overrightarrow{0}$,  suy ra $I$ là trọng tâm tam giác $ABC$. Ta có
		$$I = \left(\dfrac{3 +  2 + (-2)}{3}, \dfrac{-1 + 3 + 1}{3}, \dfrac{2 + (-3) + (-2)}{3}\right) = (1, 1, -1).$$
		 Xét biểu thức cần tối thiểu hoá\\
		$ 
		\begin{aligned}
			&\overrightarrow{MA} \cdot \overrightarrow{MB} + \overrightarrow{MB} \cdot  \overrightarrow{MC} + \overrightarrow{MC} \cdot \overrightarrow{MA} \\
			&= (\overrightarrow{MI} + \overrightarrow{IA}) \cdot (\overrightarrow{MI} + \overrightarrow{IB}) + (\overrightarrow{MI} + \overrightarrow{IB})  \cdot (\overrightarrow{MI} + \overrightarrow{IC})+ (\overrightarrow{MI} + \overrightarrow{IC}) \cdot (\overrightarrow{MI} + \overrightarrow{IA}) \\
			&= 3MI^2 + 2\overrightarrow{MI} \cdot (\overrightarrow{IA}  + \overrightarrow{IB} + \overrightarrow{IC}) + (\overrightarrow{IA} \cdot \overrightarrow{IB} + \overrightarrow{IB} \cdot \overrightarrow{IC} + \overrightarrow{IC} \cdot \overrightarrow{IA}) \\
			&= 3MI^2  + \text{ (const).}
		\end{aligned}
		$\\ 
		Vì $\overrightarrow{IA} + \overrightarrow{IB} + \overrightarrow{IC} = 0$, nên biểu thức trên đạt giá  trị nhỏ nhất khi $MI$ nhỏ nhất.\\
		Do $M$ thuộc mặt phẳng $(Oyz)$ nên $M$ là hình chiếu vuông góc của $I(1, 1, -1)$ lên mặt phẳng $(Oyz)$, suy ra  $M(0, 1, -1)$.\\
		Vậy $a - 2b + c = 0 - 2 \cdot 1  + (-1) = -3$.
	 }
\end{bt}


%%%%%%%%%%%%%%%% TỰ LUẬN %%%%%%%%%%%%%%%%

% nguồn: Lop_12_Hinh_hoc_Chuong_2
\begin{bt}%[2H2N2-2]%[ma:dad9fe57]
	Trong không gian $Oxyz$ cho điểm $A(3;-2;-1)$. Gọi $A_1$, $A_2$, $A_3$ lần lượt là hình chiếu của điểm $A$ lên các mặt phẳng  toạ độ. Tìm toạ độ các điểm $A_1$, $A_2$, $A_3$
	\loigiai{ 
		 Tọa độ của điểm $A_1=(3;-2;0)$.\\
		Toạ độ của điểm $A_2=(3;0;-1)$.\\ 
		Toạ độ của điểm  $A_3=(0;-2;-1)$.
	 }
\end{bt}

% nguồn: Lop_12_Hinh_hoc_Chuong_2
\begin{bt}%[2H2H2-2]%[ma:4c2fdccb]
	Trong không gian $Oxyz$, cho ba điểm $A(5;-3;0)$, $B(2;1;-1)$, $C(4;1;2)$. 
	\begin{enumerate}
		\item Tìm tọa độ của vectơ  $\overrightarrow{u}=2\overrightarrow{AB}+\overrightarrow{AC}-5\overrightarrow{BC}$.
		\item Tìm tọa độ điểm $N$ sao  cho $2\overrightarrow{NA}=-\overrightarrow{NB}$.
	\end{enumerate}
	 \loigiai{
		\begin{enumerate}
			\item Ta có  $\overrightarrow{AB}=(-3;4;-1)$, $\overrightarrow{AC}=(-1;4;2)$, $\overrightarrow{BC}=(2;0;3)$.\\
			Lúc đó  $2\overrightarrow{AB}=(-6;8;-2)$,$\overrightarrow{AC}=(-1;4;2)$, $-5\overrightarrow{BC}=(-10;0;-15)$.\\
			Vậy  $\overrightarrow{u}=(-6-1-10;8+4+0;-2+2-15)=(-17;12;-15)$.
			\item  Gọi $N(x;y;z)$. Ta có $\overrightarrow{NA}=(5-x;-3-y;-z)$ và $\overrightarrow{NB}=(2-x;1-y;-1-z)$.\\
			Lúc  đó
			$$2\overrightarrow{NA}=-\overrightarrow{NB} \Leftrightarrow \heva{&2\cdot  (5-x)=-(2-x)\\& 2\cdot (-3-y)=-(1-y)\\& 2\cdot (-z)=-(-1-z) }\Leftrightarrow \heva{&10-2x=x-2\\&-6-2y=y-1\\&-2z=z+1} \Leftrightarrow \heva{&3x=12\\&3y=-5\\&3z=-1} \Leftrightarrow \heva{&x=4\\&y=-\dfrac{5}{3}\\&z=-\dfrac{1}{3}.}$$
			Vậy $N\left(4;-\dfrac{5}{3};-\dfrac{1}{3}\right)$. 
		\end{enumerate}
	} 
\end{bt}

% nguồn: Lop_12_Hinh_hoc_Chuong_2
\begin{bt}%[2H2H2-2]%[ma:654fcca9]
	Trong không gian với hệ trục tọa độ $Oxyz$ cho hình hộp $ABCD.A'B'C'D'$ có  $A(1;0;1)$, $B(2;1;2)$, $D(1;-1;1)$, $C'(4;5;-5)$. Tính tọa độ đỉnh $A'$ của hình hộp.
	\loigiai{ 
		Theo quy tắc hình hộp  ta có $\overrightarrow{AB}+\overrightarrow{AD}+\overrightarrow{AA'}=\overrightarrow{AC'}\Rightarrow \overrightarrow{AA'}=\overrightarrow{AC'}-\overrightarrow{AB}-\overrightarrow{AD}$.\\
		Mặt khác  $\overrightarrow{AC'}=(3;5;-6)$, $\overrightarrow{AB}=(1;1;1)$, $\overrightarrow{AD}=(0;-1;0)$.\\
		Do đó $\overrightarrow{AA'}=(2;5;-7)$.\\ 
		Suy ra $A'(3;5;-6)$. 
	} 
\end{bt}

% nguồn: Lop_12_Hinh_hoc_Chuong_2
\begin{bt}%[2H2H2-2]%[ma:9ede1ac4]
	Một máy bay $M$ đang cất cánh từ phi trường với hệ toạ độ không gian $Oxyz$ được thiết lập như hình vẽ dưới đây. Cho  biết $M$ là vị trí của máy bay và $OM=14$, $\widehat{NOB}=32^{\circ}$, $\widehat{MOC}=65^{\circ}$. Tìm toạ độ điểm $M$.
	\begin{center}
		\begin{tikzpicture}[scale=0.8,font=\footnotesize,line join=round,line cap=round,>=stealth]
				\coordinate (O) at (0,0); 
				 \coordinate (A) at (-2,-2);
				\coordinate (B)  at (6,0);
				\coordinate (N) at  ($(B)+(A)-(O)$);
				\coordinate (C) at ($(O)+(0,4)$); 
				\coordinate (M) at  ($(C)+(N)-(O)$);
				\draw[dashed]  (A)--(N)--(M)--(C);
				\draw[dashed] (O)--(N)--(B); 
				\draw (O)--(M); 
				\foreach \i/\g in {O/110,A/90,B/90,C/180,M/0,N/-90}{\draw[fill=blue](\i) circle  (1.5pt)
					($(\i)+(\g:5mm)$) node[scale=1]{$\i$};} 
				\coordinate  (Z) at ($(O)+(0,5)$);
				 \coordinate (Y) at ($(B)+(1,0)$);
				 \coordinate (X) at ($(A)-(1,1)$);
				 \draw[->] (O)--(Z);
				\draw[->] (O)--(X); 
				\draw[->]  (O)--(Y);
				 \draw (X) node[scale=1,left]{$x$};
				\draw  (Y) node[scale=1,above]{$y$};
				 \draw (Z) node[scale=1,left]{$z$};
				\pic[draw,thin,angle  radius=2mm] {right angle = O--C--M};
				\pic[draw,thin,angle  radius=2mm] {right angle = O--A--N};
				\pic[draw,thin,angle  radius=2mm] {right angle = N--B--O};
				\pic[draw,thin,angle radius=4mm] {angle = M--O--C}; 
				\pic[draw,thin,angle  radius=6mm] {angle = N--O--B};
				\draw ($(O)+(0,1)$) node[scale=1,right]{$65^{\circ}$}; 
				\draw ($(O)+(2,0)$)  node[scale=1,below]{$32^{\circ}$};
				 \draw ($(O)!0.5!(M)$) node[scale=1,above]{$14$};
			\end{tikzpicture}
	\end{center}
	\loigiai{ 
		Ta có  $OC=OM\cdot \cos65^{\circ} \approx 5{,}9$; $ON=CM=OM\cdot \sin 65^{\circ} \approx 12{,}7$; $OB=ON\cdot \cos32^{\circ} \approx 10{,}8$.\\
		Vì $OANB$ là hình chữ nhật nên $\overrightarrow{ON}=\overrightarrow{OA}+\overrightarrow{OB}$.\\ 
		$OA=NB=ON\cdot \sin  32^{\circ}=6{,}7$\\
		Vì $OCMN$ là hình chữ nhật  nên $\overrightarrow{OM}=\overrightarrow{OC}+\overrightarrow{ON}=\overrightarrow{OA}+\overrightarrow{OB}+\overrightarrow{OC}=6{,}7\cdot \overrightarrow{i}+10{,}8\cdot \overrightarrow{j}+5{,}9\cdot \overrightarrow{k}$.\\
		Vậy $M\left(6{,}7; 10{,}8;  5{,}9\right)$.
	} 
\end{bt}

% nguồn: Lop_12_Hinh_hoc_Chuong_2
\begin{bt}%[2H2H2-2]%[ma:17a51540]
	Trong không gian $Oxyz$ cho hình hộp chữ nhật $ABCD.A'B'C'D'$ có cạnh $AB=AA'=2$, $AD=4$. Gọi $E$ là tâm của hình chữ nhật $ABCD$, $F$ là trung điểm của cạnh $AC'$. Với hệ toạ độ $Oxyz$ được thiết lập như hình dưới đây. Hãy xác định  toạ độ các đỉnh của hình hộp chữ nhật và toạ độ hai điểm $E$ và $F$?
	\begin{center}
		\begin{tikzpicture}[scale=0.8,font=\footnotesize,line join=round,line cap=round,>=stealth]
				\coordinate  (A) at (0,0);
				\coordinate (B)  at (-2,-2);
				\coordinate (D) at (6,0); 
				 \coordinate (C) at ($(B)+(D)-(A)$);
				\coordinate (E) at ($(A)!0.5!(C)$); 
				\foreach \i in {A,B,C,D}{\coordinate (\i')  at ($(\i)+(0,4)$);}
				\coordinate  (F) at ($(A')!0.5!(C)$);
				 \draw (A')--(B')--(C')--(D')--cycle;
				 \draw (B)--(B')
					 (C)--(C')
					(D)--(D') 
					 (B)--(C)--(D);
				 \draw[dashed,thin](B)--(A)--(A')
					(A)--(D); 
				\foreach \i/\g in {A'/110,B'/90,C'/90,D'/-30,A/-90,B/-90,C/-90,D/-90,E/0,F/0}{\draw[fill=blue](\i)  circle (1.5pt)
					($(\i)+(\g:5mm)$) node[scale=1]{$\i$};} 
				\coordinate  (Z) at ($(A)+(0,5)$);
				\coordinate  (Y) at ($(D)+(1,0)$);
				 \coordinate (X) at ($(B)-(1,1)$);
				\draw[->] (A')--(Z); 
				 \draw[->] (B)--(X);
				\draw[->]  (D)--(Y);
				 \draw[dashed,thin](A')--(C)--(A)--(C')
					 (B)--(D);
				\draw (X) node[scale=1,left]{$x$};  x
				\draw (Y) node[scale=1,above]{$y$}; 
				 \draw (Z) node[scale=1,left]{$z$};
				\coordinate  (I) at ($(A)-(0.4,0.4)$);
				\coordinate (J) at ($(A)+(0.5,0)$); 
				\coordinate  (K) at ($(A)+(0,0.5)$);
				\draw[->,red]  (A)--(I);
				\draw[->,red] (A)--(J); 
				\draw[->,red] (A)--(K); 
				\draw (I) node[scale=1,left]{$\vec{i}$};  x
				\draw (J)  node[scale=1,above]{$\vec{j}$};
				\draw  (K) node[scale=1,left]{$\vec{k}$};
			\end{tikzpicture}
	\end{center}
	\loigiai{ 
	Chọn $A\left(0; 0; 0\right)$ là gốc toạ độ thì toạ độ các điểm như sau $B\left(2; 0; 0\right)$; $C\left(2; 4; 0\right)$;  $D\left(0; 4; 0\right)$.\\
	$A'\left(0; 0;  2\right); B'\left(2; 0; 2\right); C'\left(2; 4; 2\right); D'\left(0; 4; 2\right).$\\
	$E$, $F$ là  trung điểm của $AC$ và $A'C$ nên $E\left(1; 2; 0\right)$ và $F\left(1; 2; 1\right)$.
	 }
\end{bt}

% nguồn: Lop_12_Hinh_hoc_Chuong_2
\begin{bt}%[2H2H2-2]%[ma:b9189a0e]
	Trong không gian $Oxyz$,  cho ba điểm $A(3;-1;2)$, $B(1;2;3)$, $C(4;-2;1)$.
	\begin{enumerate}
		\item Chứng  minh ba điểm $A$, $B$, $C$ không thẳng hàng. Xác định tọa độ trọng tâm tam giác $ABC$.
		\item Tìm tọa độ điểm $D$ biết  tứ giác $ABCD$ là hình bình hành.
		\item Tìm tọa độ giao điểm $E$ của đường thẳng $BC$ với  mặt phẳng tọa độ $(Oxz)$.
	\end{enumerate}
	\loigiai{ 
		\begin{enumerate}
			\item Ta có $\overrightarrow{AB} =(-2;3;1)$  và $\overrightarrow{AC}=(1;-1;-1)$. \\
			Vì  tỉ số $\dfrac{-2}{1} \ne \dfrac{3}{-1}$ nên hai vectơ $\overrightarrow{AB}, \overrightarrow{AC}$ không cùng phương. \\
			Vậy ba điểm $A$,  $B$, $C$ không thẳng hàng. \\
			Tọa độ trọng  tâm $G$ của tam giác $ABC$ là
			$G\left(\dfrac{3+1+4}{3};\dfrac{-1+2-2}{3};\dfrac{2+3+1}{3}\right)=G\left(\dfrac{8}{3};-\dfrac{1}{3};2\right)$. 
			\item Tứ giác $ABCD$ là hình bình hành khi và  chỉ khi $\overrightarrow{AB}=\overrightarrow{DC}$. \\
			 Gọi $D\left(x_D;y_D;z_D\right)$, ta có $\overrightarrow{DC} = \left(4-x_D;-2-y_D;1-z_D\right)$.\\
			Lúc  đó
			$$\overrightarrow{AB}=\overrightarrow{DC} \Leftrightarrow \heva{& -2=4-x_D\\&3=-2-y_D\\&1=1-z_D} \Leftrightarrow \heva{&x_D=6\\&y_D=-5\\&z_D=0 .}$$ 
			Vậy $D(6;-5;0)$. 
			\item Vì $E$ thuộc mặt phẳng  $(Oxz)$ nên $y_E=0$, suy ra $E\left(x_E;0;z_E\right)$. \\
			Vì $E$ thuộc đường thẳng $BC$ nên ba điểm $B$, $C$, $E$ thẳng hàng, do đó $\overrightarrow{BE}$ và $\overrightarrow{BC}$ cùng phương.  \\
			Ta có $\overrightarrow{BC}  = (3;-4;-2)$ và $\overrightarrow{BE}=\left(x_E-1;0-2;z_E-3\right) = \left(x_E-1;-2;z_E-3\right)$.\\
			Vì $\overrightarrow{BE}$ và $\overrightarrow{BC}$ cùng phương khi tồn tại $k$ sao cho  $\overrightarrow{BE} = k\overrightarrow{BC}$.\\
			 Lúc đó
			$$\heva{&x_E-1=3k\\&-2=-4k\\&z_E-3=-2k}  \Leftrightarrow \heva{&k=\dfrac{1}{2}\\&x_E=\dfrac{5}{2}\\&z_E=2.}$$
			Vậy $E\left(\dfrac{5}{2};0;2\right)$. 
		\end{enumerate}
	} 
\end{bt}

% nguồn: Lop_12_Hinh_hoc_Chuong_2
\begin{bt}%[2H2H2-2]%[ma:4fc162ae]
	Cho hình hộp chữ nhật $OABC.O'A'B'C'$ có cạnh $OA=4$, $OC=6$, $OO'=3$. Chọn hệ trục tọa độ $Oxyz$  có gốc tọa độ $O$; các điểm $A$, $C$, $O'$ lần lượt nằm trên các tia $Ox$, $Oy$, $Oz$. Xác định tọa độ các điểm $A$, $B$, $B'$.
	\loigiai{ 
		\begin{center}
			\begin{tikzpicture}[scale=0.7,font=\footnotesize,line join=round,
						line cap=round,>=stealth] 
					 \def\a{6}
					 \def\b{3}
					 \def\h{3.5}
					\path (0:0)  coordinate (O)
						++(0:\a)  coordinate (C)
						++(-130:\b) coordinate (B) 
						($(O)+(B)-(C)$) coordinate (A) 
						($(O)+(90:\h)$) coordinate (O') 
						($(C)+(90:\h)$)  coordinate (C')
						($(A)+(90:\h)$)  coordinate (A')
						($(B)+(90:\h)$) coordinate (B') 
						 (O')++(90:1) coordinate (z)
						(A)++(-130:0.7)  coordinate (x)
						(C)++(0:1) coordinate  (y);
					 \draw[dashed] (A)--(O)--(C)
						 (O)--(O');
					 \draw [->,dashed] (O)--(B');
					\draw (C)--(C') 
						(B)--(B') 
						 (A)--(A')
						(A)--(B)--(C) 
						 (A')--(B')--(C')--(O')--cycle;
					 \draw[->] (O')--(z);
					\draw[->]  (A)--(x);
					 \draw[->] (C)--(y);
					\foreach \x/\g in {A/180,B/0,C/-90,O/180,A'/180,B'/0,C'/0,O'/180,x/180,y/90,z/0} 
					\fill (\x)  circle (1pt)
						($(\g:4mm)+(\x)$)  node {$\x$};
				\end{tikzpicture}
		\end{center}
		Ta có 
		\begin{itemize}
			\item $\overrightarrow{OA}=4\overrightarrow{i}+0\overrightarrow{j}+0\overrightarrow{k}$ suy ra $A(4;0;0)$. 
			\item $\overrightarrow{OB}=\overrightarrow{OA}+\overrightarrow{OC}=4\overrightarrow{i}+6\overrightarrow{j}+0\overrightarrow{k}$ suy ra $B(4;6;0)$. 
			 \item $\overrightarrow{OB'}=\overrightarrow{OA}+\overrightarrow{OC}+\overrightarrow{OO'}=4\overrightarrow{i}+6\overrightarrow{j}+3\overrightarrow{k}$ suy ra $B'(4;6;3)$.
		\end{itemize}
	 }
\end{bt}

% nguồn: Lop_12_Hinh_hoc_Chuong_2
\begin{bt}%[2H2H2-2]%[ma:22f4806c]
	Trong không gian với hệ tọa độ $Oxyz$ cho hai điểm $B(1;2;-3)$, $C(7;4;-2)$. Tìm toạ độ điểm $E$ thỏa  nãm đẳng thức $\overrightarrow{CE}=2\overrightarrow{EB}$.
	 \loigiai{
		 Gọi $E(x;y;z)$.\\
		Ta có  $\overrightarrow{CE}=(x-7;y-4;z+2)$ và $2\overrightarrow{EB}=(2-2x;4-2y;-6-2z)$\\
		Khi  đó $\overrightarrow{CE}=2\overrightarrow{EB}\Leftrightarrow\heva{&x-7=2-2x\\&y-4=4-2y\\&z+2=-6-2z}\Leftrightarrow\heva{&x=3\\&y=\dfrac{8}{3}\\&z=-\dfrac{8}{3}.}$\\
		Vậy toạ độ của điểm $E$ thoả mãn hệ thức  là $E\left(3;\dfrac{8}{3};-\dfrac{8}{3}\right)$.
	} 
\end{bt}

% nguồn: Lop_12_Hinh_hoc_Chuong_2
\begin{bt}%[2H2V2-2]%[ma:fd4096cb]
	Cho hai điểm $A(3; -1; 2)$, $B(2; 3; -3)$, $C(-2; 1; -2)$ và mặt phẳng $(Oyz)$. Gọi $M(a; b; c)$ là  điểm thuộc mặt phẳng $(Oyz)$ sao cho biểu thức
	$\overrightarrow{MA} \cdot \overrightarrow{MB} + \overrightarrow{MB}  \cdot \overrightarrow{MC} + \overrightarrow{MC} \cdot \overrightarrow{MA}$ đạt giá trị nhỏ nhất. Tính $a - 2b + c$.
\end{bt}

% nguồn: Lop_12_Hinh_hoc_Chuong_2
\begin{bt}%[2H2V2-2]%[ma:176c6519]
	Trong không gian $Oxyz$, cho tam  giác $ABC$ với $A(1;0;0)$, $B(0;0;1)$ và $C(2;1;1)$. Tìm tâm đường tròn ngoại tiếp tam giác $ABC$.
	 \loigiai{
		Ta có $\overrightarrow{AB} =  (-1; 0; 1)$, $\overrightarrow{AC} = (1; 1; 1)$. \\
		Mặt phẳng  $(ABC)$ có vectơ pháp tuyến
		$\overrightarrow{n} = \left[\overrightarrow{AB},  \overrightarrow{AC}\right]=(-1; 2; -1)$. \\
		Mặt phẳng $(ABC)\colon -1(x-1) + 2y - z = 0 \Leftrightarrow -x  + 2y - z + 1 = 0$. \\
		Gọi $I$  là tâm đường tròn ngoại tiếp tam giác $ABC$.\\
		 Ta có
		 \allowdisplaybreaks
		\begin{eqnarray*}
			&& \heva{&IA = IB \\ &IB = IC \\ &I \in (ABC)}\\ &\Leftrightarrow &\heva{&(a-1)^2 + b^2 + c^2 = a^2 + b^2 + (c-1)^2 \\  &a^2 + b^2 + (c-1)^2 = (a-2)^2 + (b-1)^2 + (c-1)^2 \\ &-a + 2b - c + 1 = 0}\\
			&\Leftrightarrow& \heva{&-2a + 2c  = 0 \\ &4a + 2b = 5 \\ &-a + 2b - c + 1 = 0}\\
			&\Leftrightarrow& \heva{&a = c \\ &4a + 2b = 5 \\ &-2a + 2b = -1}\\ 
			&\Leftrightarrow& \heva{&a = 1 \\ &b = \dfrac{1}{2} \\ &c =  1.}
		\end{eqnarray*}
		 Vậy $I\left(1; \dfrac{1}{2}; 1\right)$.
	 }
\end{bt}
